% Portfolio Sentinel -- Design Document
% Updated at the end of every milestone / sub-milestone (see CLAUDE.md).
% Compile: pdflatex design_doc.tex (run twice for the table of contents), or upload to Overleaf.

\documentclass[11pt]{article}

\usepackage[T1]{fontenc}
\usepackage[utf8]{inputenc}
\usepackage{lmodern}
\usepackage[margin=1in]{geometry}
\usepackage{amsmath,amssymb}
\usepackage{booktabs}
\usepackage{tabularx}
\usepackage{enumitem}
\usepackage{xcolor}
\usepackage{listings}
\usepackage[most]{tcolorbox}
\usepackage[hidelinks]{hyperref}

\setlength{\parskip}{0.5em}
\setlength{\parindent}{0pt}
\setlist{itemsep=0.2em}

% ---------- helpers ----------
\newcommand{\code}[1]{\texttt{\detokenize{#1}}}

\newtcolorbox{keyidea}[1][Key idea]{colback=blue!4,colframe=blue!50!black,
  fonttitle=\bfseries,title=#1,breakable}
\newtcolorbox{whyitmatters}[1][Why it matters]{colback=green!4,colframe=green!40!black,
  fonttitle=\bfseries,title=#1,breakable}
\newtcolorbox{takeaway}[1][Key takeaway]{colback=orange!5,colframe=orange!70!black,
  fonttitle=\bfseries,title=#1,breakable}
\newtcolorbox{reviewq}[1][Review questions]{colback=gray!6,colframe=gray!60!black,
  fonttitle=\bfseries,title=#1,breakable}

\lstdefinestyle{base}{
  basicstyle=\ttfamily\small,
  columns=fullflexible,
  keepspaces=true,
  breaklines=true,
  frame=single,
  framerule=0.3pt,
  rulecolor=\color{gray!50},
  backgroundcolor=\color{gray!5},
  showstringspaces=false,
  commentstyle=\color{green!40!black},
  keywordstyle=\color{blue!60!black}\bfseries,
  stringstyle=\color{red!50!black},
  xleftmargin=0.5em,
  xrightmargin=0.5em,
}
\lstdefinestyle{py}{style=base,language=Python}
\lstdefinestyle{plain}{style=base}

\title{\textbf{Portfolio Sentinel}\\[0.3em]\large Design Document: concepts, architecture, and decisions}
\author{Aditya Kalkar}
\date{Last updated: 2026-10-24 (through milestone M3)}

\begin{document}
\maketitle

\begin{abstract}
This document explains everything that goes into Portfolio Sentinel: the finance, the code,
the AI and agent ideas, MCP, Git/GitHub, and why each decision was made. It is written to be
read cold by someone new to the project. Part~I is background reference that grows as new
concepts appear. Part~II is a milestone-by-milestone log of what was built, how it works, and
how it was verified. Each milestone ends with review questions.
\end{abstract}

\tableofcontents
\newpage

% =====================================================================
\section{The big picture}
% =====================================================================

\subsection{What we are building}

Portfolio Sentinel is an \textbf{AI agent for portfolio managers}. A portfolio manager (PM)
runs an investment fund. They might ask:

\begin{itemize}
  \item ``Where are we most overweight versus the benchmark?''
  \item ``Are we breaching any risk limits right now?''
  \item ``If I sell \$2M of NVDA, do we still breach anything?''
\end{itemize}

The agent answers by calling \textbf{tools} that compute the answer in ordinary Python code.
Every weekday morning it also runs on its own and writes a \textbf{morning brief}: a short
report on what changed overnight and what needs attention.

\subsection{Why this project exists}

The project demonstrates four engineering capabilities:

\begin{enumerate}
  \item \textbf{Agents}: a program where a large language model (LLM) decides which tools to
        call, in what order, and explains the results.
  \item \textbf{MCP} (Model Context Protocol): an open standard for exposing tools and data
        to AI models. We build our own MCP server.
  \item \textbf{Automation}: tests on every push (CI) and a scheduled morning brief, both
        using GitHub Actions.
  \item \textbf{Metrics / evals}: measuring how good the agent actually is (accuracy, tool
        choice, made-up numbers, cost, speed) and improving it with evidence.
\end{enumerate}

The finance is deliberately simple arithmetic. The hard and interesting part is the agent
engineering, and proving it works.

\begin{keyidea}[The one rule: the agent never does math]
LLMs are good at language and at deciding what to do next. They are unreliable at
arithmetic, and they sometimes produce plausible numbers that are wrong (``hallucination'').
So every number the agent reports must come from a tool result, computed by tested Python
code. The model's job is to choose tools and explain what they return. Most of the project's
architecture follows from this rule.
\end{keyidea}

\subsection{Architecture at a glance}

\begin{lstlisting}[style=plain]
  PM question / daily schedule
              |
              v
   +----------------------+        +----------------------+
   |  Agent (harness)     |<------>|  Claude API          |
   |  agent loop, prompts,|        |  (model picks tools, |
   |  logging, limits     |        |   writes answers)    |
   +----------+-----------+        +----------------------+
              | MCP (stdio)
              v
   +----------------------+
   |  MCP server          |  thin layer: load data, call analytics, format output
   +----------+-----------+
              v
   +----------------------+        +----------------------+
   |  analytics.py        |<-------|  data.py             |
   |  pure functions,     |        |  portfolios, prices, |
   |  all the math        |        |  benchmarks, sectors |
   +----------------------+        +----------------------+
\end{lstlisting}

Read it bottom-up:

\begin{description}[style=nextline]
  \item[\code{data.py} (built in M1a)] The only code that reads files. It loads the
        portfolio, benchmark, securities list, and prices.
  \item[\code{analytics.py} (built in M1b)] Pure math functions: weights, exposure, benchmark
        comparison, limit checks, what-if trades. No files, no network, no AI.
  \item[MCP server (built in M2a)] Wraps analytics as tools an AI can call. It converts decimals
        (0.142) to percentages (14.2) and rounds. It does no new math.
  \item[Agent (built in M3)] A loop we write by hand: send the question and tool list to Claude;
        if Claude asks for a tool, run it through MCP and send back the result; repeat until
        Claude gives a final answer.
  \item[Evals (M4) and morning brief (M5)] Built on top of the agent.
\end{description}

\begin{whyitmatters}[Why layers?]
Each layer has one job and strict rules, so each can be tested on its own. If a number is
wrong, you know where to look: the math lives in exactly one place (\code{analytics.py}),
file reading lives in exactly one place (\code{data.py}), and the AI never touches either.
This is called \emph{separation of concerns}.
\end{whyitmatters}

\subsection{Milestone map}

\begin{center}
\begin{tabularx}{\textwidth}{@{}l X l@{}}
\toprule
\textbf{Milestone} & \textbf{What it delivers} & \textbf{Status} \\
\midrule
Setup & Repo created; spec reviewed; 8 decisions agreed & Done \\
M0  & Project skeleton (packaging, folders) and CI on GitHub & Done \\
M1a & Synthetic sample data with a planted story; data loaders & Done \\
M1b & Analytics core, written test-first; sample data sanity tests & Done \\
M2a & MCP server with 8 tools (task 2.1), tested in-process and over stdio & Done \\
M2b & Connected to Claude Code; tool description pass (tasks 2.2, 2.3) & Done \\
M3  & Hand-written agent loop, system prompt, run logging & Done \\
M4  & Eval harness: test cases, scoring, results history & Next \\
M5  & Morning brief, scenario evals, scheduling, delivery & \\
M6  & Real data (ETF holdings, SEC 13F filings, prices) & \\
M7  & Streamlit UI, metrics dashboard, deploy, demo & \\
M8  & Stretch goals (model comparison, prompt-injection tests, ...) & \\
\bottomrule
\end{tabularx}
\end{center}

% =====================================================================
\part{Background reference}
% =====================================================================

\section{Finance primer}

This section covers every finance idea the project uses. Nothing here is beyond basic
arithmetic; the words are the hard part.

\subsection{Funds, portfolio managers, and positions}

A \textbf{fund} pools money from investors and invests it. The \textbf{portfolio manager}
decides what the fund owns. Everything the fund owns at one moment is its
\textbf{portfolio}, and each holding is a \textbf{position}.

A position is described by:
\begin{itemize}
  \item \textbf{Ticker}: the short stock symbol (NVDA = NVIDIA, UNH = UnitedHealth).
  \item \textbf{Shares}: how many units of the stock the fund holds.
  \item \textbf{Price}: the stock's price. We use the \textbf{close}, the last traded price
        of the day.
  \item \textbf{Market value} $=$ shares $\times$ price, in dollars.
\end{itemize}

Funds also hold some \textbf{cash}, uninvested money that is ready for trades or for
investors who want their money back.

\subsection{NAV and weights}

\begin{keyidea}[Definitions]
\textbf{NAV} (net asset value) $=$ sum of all position market values $+$ cash.\\
\textbf{Weight} of a position $=$ market value $/$ NAV.
\end{keyidea}

Weight answers ``what share of the fund is this?'' In our sample fund on 2026-10-21:

\begin{itemize}
  \item NAV $=$ \$50{,}000{,}000.
  \item NVDA: 38{,}889 shares $\times$ \$180.00 $=$ \$7{,}000{,}020.
  \item NVDA weight $= 7{,}000{,}020 / 50{,}000{,}000 = 0.1400 = 14.00\%$.
\end{itemize}

\textbf{Convention used in this project.} Inside the code, weights are decimals (0.14).
They become percentages (14.0) only at the edge where the AI sees them, in fields whose names
end in \code{_pct}. Money fields end in \code{_usd}. Using one representation everywhere
inside the code prevents the classic bug of mixing 0.14 and 14 in the same calculation.

\textbf{Weights drift.} If NVDA's price rises 10\% and nothing else moves, NVDA's market value
rises, NAV rises by less (NVDA is only part of the fund), and NVDA's weight goes up even
though nobody traded. That is exactly what happens in our planted story on 2026-10-22.

\subsection{Long and short positions}

\begin{itemize}
  \item \textbf{Long}: you own the stock. You profit if it rises. Positive shares.
  \item \textbf{Short}: you borrow shares, sell them now, and buy them back later. You profit
        if the price \emph{falls}. In this project a short has negative shares, negative
        market value, and negative weight.
\end{itemize}

Funds often short a stock as a \textbf{hedge}, to offset some risk. Our sample fund is
short INTC (Intel) at about $-2\%$ of NAV. It is long several chip stocks (NVDA, AMD,
AVGO), so a short on another chip company offsets some of that bet. On 2026-10-22, INTC
falls 3\%, so the short \emph{makes} money.

\textbf{Shorts and cash.} Selling borrowed shares brings cash in. This is why NAV
$=$ market values $+$ cash still works with shorts: the short's negative market value is the
cost of buying the shares back.

\subsection{Exposure: long, short, gross, net}

Exposure measures how much of the fund's money is at risk in the market.

\begin{center}
\begin{tabular}{@{}ll@{}}
\toprule
\textbf{Measure} & \textbf{Definition} \\
\midrule
Long exposure  & sum of positive weights \\
Short exposure & sum of negative weights (a negative number) \\
Gross exposure & long $-$ short \quad (total size of all bets) \\
Net exposure   & long $+$ short \quad (direction: how much we lean on the market) \\
\bottomrule
\end{tabular}
\end{center}

\textbf{Worked example} (this is a test case in M1b): \$100 long, \$20 short, \$20 cash.
\begin{itemize}
  \item NAV $= 100 - 20 + 20 = 100$.
  \item Long $= 100\%$, short $= -20\%$.
  \item Gross $= 100 - (-20) = 120\%$; net $= 100 + (-20) = 80\%$.
\end{itemize}

Our sample fund on 2026-10-21: long 98\%, short $-2\%$, cash 4\%, so gross $= 100\%$ and
net $= 96\%$.

\subsection{Benchmarks and active weight}

Most funds are judged against a \textbf{benchmark}, usually a stock market
\textbf{index} such as the S\&P~500 (about 500 large US companies, each weighted by its size).
The question is not just ``did we make money?'' but ``did we beat the benchmark?''

\begin{keyidea}[Active weight]
Active weight $=$ portfolio weight $-$ benchmark weight.\\
Positive $=$ \textbf{overweight} (we hold more than the index). Negative $=$
\textbf{underweight}.
\end{keyidea}

Examples from our data (2026-10-21):
\begin{itemize}
  \item NVDA: fund 14.0\%, benchmark 12.0\%, active $+2.0\%$ (overweight).
  \item AAPL: fund 8.5\%, benchmark 11.0\%, active $-2.5\%$ (underweight).
  \item AMD: fund 2.5\%, not in the benchmark (0\%), active $+2.5\%$.
  \item ORCL: not held (0\%), benchmark 2.0\%, active $-2.0\%$. Not owning something the index
        owns is also a bet.
\end{itemize}

Active weights are the PM's \textbf{bets}. A big active weight means the fund's results can
differ a lot from the index's. That is good if the bet is right and painful if it is wrong.

\subsection{Sectors}

Companies are grouped into \textbf{sectors} by industry. The standard scheme is
\textbf{GICS} (Global Industry Classification Standard), which has 11 sectors. Our data uses
8 of them: Information Technology, Communication Services, Consumer Discretionary, Health
Care, Financials, Consumer Staples, Industrials, Energy.

Sector exposure adds up weights by sector. Sector active weight compares that total with the
benchmark's sector total. Stocks in the same sector tend to move together, so a big sector bet
is a concentrated risk even when each single stock looks small.

\subsection{Risk limits, warnings, and breaches}

Firms set \textbf{risk limits} so one bad bet cannot sink the fund. Ours (decimals in the
portfolio file):

\begin{center}
\begin{tabularx}{\textwidth}{@{}l X l@{}}
\toprule
\textbf{Limit} & \textbf{Meaning} & \textbf{Value} \\
\midrule
\code{max_position}      & no single position larger than this (by absolute weight) & 15\% \\
\code{max_stock_active}  & no stock more than this over- or underweight vs benchmark & 5\% \\
\code{max_sector_active} & no sector more than this over- or underweight & 8\% \\
\code{max_gross}         & gross exposure at most this & 115\% \\
\code{min_net}/\code{max_net} & net exposure must stay in this range & 90\%--102\% \\
\bottomrule
\end{tabularx}
\end{center}

A \textbf{breach} means the limit is exceeded. A \textbf{warning} means it is close, so the
PM can act before it becomes a breach:

\begin{itemize}
  \item For the \code{max_*} limits: warn at 90\% of the limit or more (\code{warn_at} $= 0.9$).
        For a 15\% position limit, the warning line is $0.9 \times 15\% = 13.5\%$.
  \item For the net range: warn within 2 percentage points of either end
        (\code{net_warn_buffer} $= 0.02$), so below 92\% or above 100\%.
\end{itemize}

\textbf{Why net uses a fixed buffer.} The net range is narrow (90--102\%). ``90\% of the
limit'' would put the warning line for \code{max_net} at $0.9 \times 102\% = 91.8\%$, so
nearly every normal fund would always be ``in warning.'' A fixed 2-point distance gives
sensible warnings.

\textbf{Flag identity} (a decision from the spec review). A flag is identified by its rule and its
subject, for example (\code{max_position}, NVDA). If a trade turns a breach into a warning on
the same rule and stock, that counts as ``breach resolved,'' not a new flag.

\textbf{Size, not sign} (added in M1b). The \code{max_*} limits are about how \emph{big} a bet
is, whichever direction it points:
\begin{itemize}
  \item A short at $-16\%$ of NAV breaches a 15\% position limit just as a $+16\%$ long does.
  \item Being $5.5\%$ \emph{underweight} a stock breaches a 5\% active limit just as being
        $5.5\%$ overweight does. Not owning a stock the index owns is a bet too.
\end{itemize}
So the checks compare the absolute value, $|\text{value}|$, with the limit, but each flag
reports the signed value, so a reader can tell a $-14\%$ short from a $+14\%$ long.

\textbf{Exactly at the limit is a warning, not a breach.} ``Max 15\%'' means 15\% is allowed.
A breach is strictly above the limit; a warning is at or above the warning line. For
\code{min_net}, the direction flips: a breach is strictly below 90\%.

\subsection{What-if trades}

A \textbf{what-if} (or pre-trade) analysis asks: if I made this trade, what would the
portfolio look like? In this project trades are in dollars. Positive means buy (or cover a
short). Negative means sell (or go short). Cash moves the opposite way. A \$1M sale leaves NAV
unchanged, because \$1M of stock becomes \$1M of cash, but it lowers that stock's weight.

Simplifications we state openly: no transaction costs, no taxes, no settlement delay, and
trades happen at the closing price.

\textbf{Dollar trades mean fractional shares.} Selling \$1M of NVDA at \$198.00 is
$1{,}000{,}000 / 198 = 5{,}050.505\ldots$ shares. Real trades are in whole shares, but for a
``what would happen'' question the difference (under \$200 here) does not change the answer,
and dollars are how PMs naturally phrase trades.

In the industry this is called a \textbf{pre-trade compliance check}: before an order is sent,
the system simulates it and checks the result against the fund's limits. Our
\code{simulate_trades} tool (M2) is a small version of exactly that.

\subsection{Returns}

\textbf{Price return} over a period $=$ end price $/$ start price $- 1$. For NVDA from
\$180.00 to \$198.00: $198/180 - 1 = 0.10 = +10\%$. The fund's own return is roughly the
weighted sum of its positions' returns. NVDA at 14\% weight rising 10\% adds about
$0.14 \times 10\% = 1.4\%$ to the fund.

% ---------------------------------------------------------------------
\section{AI, agents, and MCP primer}

These ideas are introduced briefly here and covered in depth in M2 (MCP) and M3 (agents).

\subsection{Large language models}

An \textbf{LLM} (large language model) such as Claude is trained to predict text. Given a
conversation, it produces the next message. It does not ``look up'' facts or run code by
itself; it generates likely text. That is why it can write fluent explanations, and also why
it can confidently produce a wrong number.

\textbf{Tokens.} Models read and write text in small pieces called tokens (roughly
three-quarters of a word). API usage is billed per million input and output tokens.

\subsection{Tool use}

\textbf{Tool use} (also called function calling) lets a model ask your program to run a
function. You describe each tool (its name, what it does, and its inputs as a JSON schema).
The model replies either with text or with a request such as ``call \code{check_limits} with
\code{portfolio="sample_fund"}''. Your code runs the function and sends the result back. The
model never runs code itself; your program stays in control.

\subsection{Agents and the agent loop}

An \textbf{agent} is an LLM in a loop with tools:

\begin{enumerate}
  \item Send the conversation and the tool descriptions to the model.
  \item If the model asks for a tool, run it and add the result to the conversation.
  \item Repeat until the model gives a final answer, or a step limit is reached.
\end{enumerate}

The program around the model (the loop, the logging, the safety limits) is called the
\textbf{harness}. We write ours by hand, with no framework such as LangChain, so that every line
is understood and explainable.

\subsection{The Messages API, block by block (added in M3)}

Our agent talks to Claude through one HTTP endpoint, the \textbf{Messages API}, via
Anthropic's Python SDK. Each request carries:

\begin{description}[style=nextline]
  \item[\code{model}] Which model, e.g. \code{claude-opus-5-5}.
  \item[\code{system}] The \textbf{system prompt}: standing instructions that apply to the
        whole conversation (our rules: never compute numbers, no advice, \ldots).
  \item[\code{tools}] The tool definitions: name, description, input schema.
  \item[\code{messages}] The whole conversation so far, alternating \code{user} and
        \code{assistant}. The API is \textbf{stateless}: it remembers nothing between
        requests, so the full history is sent every time.
  \item[\code{max_tokens}] A ceiling on the length of the reply.
\end{description}

A message's content is a list of \textbf{content blocks}, each with a \code{type}:

\begin{center}
\begin{tabular}{@{}lll@{}}
\toprule
\textbf{Block type} & \textbf{Who writes it} & \textbf{What it is} \\
\midrule
\code{text}        & either    & ordinary text \\
\code{tool_use}    & assistant & ``please run tool X with these arguments''; has an \code{id} \\
\code{tool_result} & user (our code) & the tool's output, tagged with the same \code{id} \\
\code{thinking}    & assistant & the model's private reasoning (we just pass it back) \\
\bottomrule
\end{tabular}
\end{center}

Every response also has a \textbf{\code{stop_reason}}: why the model stopped writing.
\code{tool_use} means ``I'm waiting for tool results''; \code{end_turn} means ``that's my
answer''. Others: \code{max_tokens} (hit the length ceiling), \code{refusal} (declined). The
whole agent loop is driven by this one field.

\textbf{Thinking and effort.} Claude Opus 5.5 can reason privately before answering; how
much is set by \code{effort} (\code{low} to \code{max}). Higher effort is slower and costs
more tokens. We use \code{medium}.

\textbf{Tokens and cost.} The response reports \code{usage.input_tokens} (everything sent:
system prompt, tool definitions, history) and \code{usage.output_tokens} (what the model
wrote). Claude Opus 5.5 costs \$4 per million input tokens and \$20 per million output tokens.
Because the API is stateless, every step of the loop pays again for the system prompt and
tool definitions. \textbf{Prompt caching} (a later optimisation) lets repeated prefixes be
billed at a fraction of the price.

\textbf{Refusals and fallbacks.} Safety classifiers can occasionally decline a request, even a
harmless one. With \code{fallbacks="default"} (a beta feature), the API quietly retries such
a request on another model and returns that answer instead of nothing.

\subsection{MCP (Model Context Protocol)}

\textbf{MCP} is an open standard, created by Anthropic, for connecting AI applications to tools
and data. It is often compared to USB-C: one standard plug instead of a custom integration for
every pair of app and tool.

\begin{itemize}
  \item An \textbf{MCP server} offers tools (and optionally data and prompts). Ours will offer
        8 portfolio tools.
  \item An \textbf{MCP client} connects to servers and lets a model use their tools. Claude
        Code is a client, and our agent will be one too.
  \item \textbf{stdio transport}: the client starts the server as a subprocess and they talk
        over standard input/output. Simple and local.
\end{itemize}

Because we use the standard, the same server works in Claude Code (M2) and in our own agent
(M3) with no changes.

\subsection{Inside an MCP tool (added in M2a)}

When a client connects, it asks the server for its tool list (\code{list_tools}). For each
tool, the server sends back four things:

\begin{description}[style=nextline]
  \item[Name] e.g. \code{check_limits}. The model uses it to ask for the tool.
  \item[Description] Plain text explaining what the tool is for. This is the \emph{only}
        thing the model reads when deciding which tool fits a question, so a vague
        description means wrong tool choices. Ours always end with an example question.
  \item[Input schema] A \textbf{JSON Schema}: a standard JSON description of the arguments,
        their types and which are required, e.g. ``\code{trades} is an array of objects, each
        with a string \code{ticker} and a number \code{dollars}''. The SDK builds it
        automatically from the Python type hints.
  \item[Annotations] Hints about behaviour. \code{read_only_hint=True} tells the client
        the tool changes nothing, so it is safe to call without asking the user.
\end{description}

The server can also send \textbf{instructions} once, at connection time: general guidance
for the model that applies to every tool. Ours explain the units (\code{_pct} is percent of
NAV, \code{_usd} is dollars, negative means short) and say to report numbers exactly as given.

\textbf{Structured output.} A tool can return plain text, or \textbf{structured content}: a
JSON object with named fields, such as \code{{"value_pct": 15.2, "status": "breach"}}.
Named fields are easier for a model to read correctly than prose, and easy for tests to
check. In MCP SDK v2, a tool gets structured output by declaring its return type as
\code{dict[str, Any]}.

\textbf{Errors.} If a tool raises \code{ToolError("No prices for 2026-10-17. Available
dates: ...")}, the client receives that message marked as an error, and the model can read it
and retry. Any \emph{other} exception is hidden: the client only sees ``Error executing tool
check\_limits''. This is a security default (internal error messages can leak details), so we
convert our own \code{ValueError}s, whose messages are written for the model, into
\code{ToolError}s on purpose.

\textbf{Testing a server.} The SDK's \code{Client} can connect to a server object directly
in the same Python process (\emph{in-process}), with no subprocess and no pipes. Tests run
fast and go through the real protocol. One test also starts the server as a subprocess over
stdio, which is exactly how Claude Code will run it.

\subsection{Connecting a client: configuration and trust (added in M2b)}

A client needs to know \emph{how to start} a stdio server: which command to run, with which
arguments. Claude Code stores this in one of three \textbf{scopes}:

\begin{description}[style=nextline]
  \item[Local] Only you, only this project (stored in your user settings).
  \item[Project] Everyone who works on this repo: a file called \code{.mcp.json} in the repo
        root, committed with the code. This is what we use.
  \item[User] Only you, but in every project on your machine.
\end{description}

Our \code{.mcp.json}:

\begin{lstlisting}[style=plain]
{"mcpServers": {"pm-agent": {"type": "stdio",
                             "command": ".venv/bin/python",
                             "args": ["-m", "pm_agent.mcp_server"]}}}
\end{lstlisting}

The command is a \emph{relative} path, so it works for anyone who clones the repo and creates
the virtual environment in the usual place.

\textbf{Trust.} An MCP server is a program running on your machine, so it can do anything
you can do. Claude Code therefore asks before starting a project server for the first time,
and asks again before each tool call unless you allow it. The prompt offers ``use this and all
future MCP servers in this project''. We did not choose that: the repo is public, so a pull
request could add a server to \code{.mcp.json}, and auto-approval would start it without
asking. Approving servers one at a time is safer.

\subsection{Evals}

An \textbf{eval} is a test suite for AI behaviour: a set of questions with known correct
answers, run against the agent and scored. Unlike unit tests, the system is not
deterministic, so we measure rates: accuracy, how often the agent picks the right tool first,
how often it invents numbers, cost, and latency. Evals turn ``it seems better'' into ``accuracy
went from 72\% to 88\%.''

\subsection{Claude Pro versus the Anthropic API}

These are separate products. Claude Pro is a subscription for using Claude yourself (claude.ai
and Claude Code), and it covers building this project. The \textbf{Anthropic API} is
pay-per-token access for your own programs, set up in the Claude Console with its own billing
and an API key. The agent we build from M3 onward runs as its own program, so it needs an API
key. The key lives in a \code{.env} file that Git ignores, so it is never committed.

% ---------------------------------------------------------------------
\section{Software engineering primer}

\subsection{Python packages and the src layout}

A \textbf{package} is a folder of Python modules you can \code{import}. Ours is
\code{pm_agent}, in \code{src/pm_agent/}. Putting the package under \code{src/} (the
\emph{src layout}) means tests can only import the package the way users would, after it is
installed, not by accident from the current folder. That catches packaging mistakes early.

\code{__init__.py} marks a folder as a package. Ours holds only a docstring.

\subsection{Virtual environments}

A \textbf{virtual environment} (\code{.venv/}) is a private copy of Python with its own
installed libraries, separate from the system Python and from other projects. This avoids
``project A needs pandas 2, project B needs pandas 3'' conflicts. Create one with
\code{python -m venv .venv}. Activate it with \code{source .venv/bin/activate}. It is
listed in \code{.gitignore}, because each machine builds its own.

\subsection{pyproject.toml and dependencies}

\code{pyproject.toml} is the standard file that describes a Python project: its name,
version, required Python version, and \textbf{dependencies} (libraries it needs).
\textbf{Version specifiers} constrain which versions are acceptable:
\begin{itemize}
  \item \code{mcp>=2.3,<3}: version 2.3 or newer, but not 3.x. A major version (the first
        number) may contain breaking changes, so we stay below 3 on purpose.
  \item \code{pandas>=2.2}: anything from 2.2 up.
\end{itemize}

\textbf{Optional dependencies} (``extras'') are groups you install only when needed. Our
\code{dev} extra adds \code{pytest}, which developers need but users of the package do not.

\code{pip install -e ".[dev]"} installs the project in \textbf{editable} mode (\code{-e}),
which links to your source folder so code changes take effect without reinstalling, and
includes the \code{dev} extra.

\subsection{Testing with pytest}

A \textbf{unit test} is a small function that checks one behaviour. \textbf{pytest} finds
files named \code{test_*.py}, runs every function named \code{test_*}, and reports passes
and failures. Useful features we use:

\begin{itemize}
  \item \code{assert x == y}: the basic check. pytest shows both values when it fails.
  \item \code{pytest.approx(1.0)}: compare decimals with a small tolerance, since
        $0.1 + 0.2$ is not exactly $0.3$ in floating point.
  \item \code{pytest.raises(ValueError, match="...")}: check that code raises an error whose
        message matches a pattern (a regular expression; \code{.*} means ``anything'').
  \item \textbf{Fixtures}, such as \code{tmp_path} (a fresh temporary folder) and
        \code{monkeypatch} (temporarily change environment variables or attributes, undone
        after the test).
\end{itemize}

\textbf{Known-answer tests} use inputs small enough to check by hand (two \$50 positions plus
\$100 cash gives 25\% each). If the test passes, you know the code is right, not just
consistent with itself.

\textbf{Test-first (TDD).} Write the test, watch it fail, then write the code. Watching it fail
proves the test actually checks something. A test that can never fail is worthless.

\textbf{Parametrized tests} (added in M1b). \code{@pytest.mark.parametrize} runs one test
function several times with different inputs, each reported as its own test:

\begin{lstlisting}[style=py]
@pytest.mark.parametrize("a_dollars, cash, expected", [
    (89, 11, [("min_net", "breach")]),    # 89% < 90%
    (91, 9,  [("min_net", "warning")]),   # within 2 points of 90%
    (95, 5,  []),                         # comfortably inside
])
def test_limits_net_uses_fixed_buffer(a_dollars, cash, expected): ...
\end{lstlisting}

This is the natural shape for boundary rules: one row per side of each line.

\textbf{Testing the tests (a mutation check).} Some tests pass the first time they run, for
example sanity checks on data that is already correct. To prove such a test can fail,
deliberately break the thing it guards (a \emph{mutation}), run it, watch it fail, and undo the
break. In M1b we changed NVDA's 10/22 price in the data file, saw three sanity tests fail, and
restored the file with \code{git checkout}.

\subsection{Dataclasses and immutability}

A \textbf{dataclass} is a Python class whose job is to hold data. The \code{@dataclass}
decorator writes the boilerplate (the constructor, a readable printout, equality) from a list
of typed fields:

\begin{lstlisting}[style=py]
@dataclass(frozen=True)
class Portfolio:
    name: str
    positions: dict[str, float]   # ticker -> shares
    cash: float
    benchmark: str | None = None  # fields with defaults come last
\end{lstlisting}

\code{frozen=True} makes instances \textbf{immutable}: \code{port.cash = 0} raises an error.
To get a changed copy you call \code{dataclasses.replace(port, cash=0)}, which builds a new
object and leaves the old one alone. Immutability matters for what-if trades: the
``before'' portfolio must stay exactly as it was so it can be compared with the ``after'' one.

One catch: \code{frozen} stops you reassigning a field, but a \code{dict} inside is still
mutable (\code{port.positions["NVDA"] = 0} would work). So code that changes positions must
copy the dict first (\code{dict(port.positions)}), and \code{apply_trades} does.

\subsection{pandas essentials used in the analytics}

\begin{description}[style=nextline]
  \item[Series] A one-column table with labels (the \textbf{index}). Prices are a Series
        indexed by ticker, so \code{prices["NVDA"]} is NVDA's price.
  \item[Index alignment] When you combine Series, pandas lines rows up by label, not by
        position. \code{pd.DataFrame({"portfolio": held, "benchmark": bench})} produces one
        row for every ticker in \emph{either} Series, with \code{NaN} (``not a number'',
        meaning missing) where one side has no value. \code{.fillna(0.0)} turns those gaps
        into zero weights. That one line handles both ``held but not in the benchmark'' and
        ``in the benchmark but not held''.
  \item[groupby] \code{w.groupby("sector")["long"].sum()} splits rows into groups by sector
        and adds up each group, like a pivot table.
  \item[clip] \code{weights.clip(lower=0)} replaces negatives with 0 (keeps only longs);
        \code{clip(upper=0)} keeps only shorts.
  \item[reindex] \code{table.reindex(new_labels, fill_value=0.0)} gives a table exactly
        those rows, adding any missing ones filled with zeros.
  \item[Stable sort] \code{sort_values(..., kind="stable")} keeps rows that tie in their
        original order. Python's \code{sorted()} is always stable. We rely on this to list
        breaches before warnings without scrambling the order within each group.
\end{description}

\subsection{Floating-point numbers}

Computers store decimals in binary, so most are slightly off: \code{0.1 + 0.2} is
\code{0.30000000000000004}. Two consequences in this project:

\begin{itemize}
  \item \textbf{Tests compare with a tolerance}: \code{pytest.approx}.
  \item \textbf{``Is it zero?'' needs a tolerance too.} Selling a whole position computes
        $\text{shares} - \text{dollars}/\text{price}$, which can leave a crumb such as
        $10^{-13}$ shares instead of exactly 0. \code{math.isclose(shares, 0, abs_tol=1e-9)}
        treats anything that small as zero, so the position is removed.
\end{itemize}

Limit checks do \emph{not} use a tolerance: a weight of 15.0000001\% is over 15\%. In
practice, the data never sits that close to a line, which is what the ``clear room'' spec
decision guarantees.

\subsection{Decorators, context managers, TypedDict, and asyncio (added in M2a)}

\begin{description}[style=nextline]
  \item[Decorator] A function that takes a function and returns it (possibly registered or
        wrapped). Writing \code{@tool} above \code{def check_limits(...)} is shorthand for
        \code{check_limits = tool(check_limits)}. Our \code{tool} decorator registers the
        function with the MCP server and returns it unchanged.
  \item[Context manager] An object used with \code{with}, which runs setup before a block
        and cleanup (or error handling) after it, even if the block raises. Ours,
        \code{readable_errors()}, catches a \code{ValueError} raised anywhere inside the
        block and re-raises it as a \code{ToolError}. \code{@contextmanager} lets you write
        one as a generator: code before \code{yield} runs on entry, and a \code{try}
        around the \code{yield} sees errors from the block.
  \item[TypedDict] A type hint for a dict with known keys, e.g. a trade must have a
        \code{ticker} (string) and \code{dollars} (number). At runtime it is an ordinary
        dict; the benefit is that tools (here, the MCP SDK building a JSON Schema) can read
        the expected keys. It is part of Python's standard library, so it adds no dependency.
  \item[async / await and asyncio] MCP clients are \emph{asynchronous}: functions declared
        \code{async def} can pause while waiting for messages (\code{await}) and let other
        work run. \code{asyncio.run(main())} starts an event loop, runs \code{main} to the
        end, and returns its result. Our tests wrap each client call in \code{asyncio.run},
        so the test functions themselves stay ordinary (no async test plugin needed).
\end{description}

\subsection{Fakes, environment variables, and JSON Lines (added in M3)}

\begin{description}[style=nextline]
  \item[Test double (fake)] A stand-in for something slow, costly, or unpredictable, used in
        tests. Our \code{FakeModel} plays Claude's part from a script (``first ask for
        \code{check_limits}, then answer''), so the agent loop can be tested in a second, for
        free, with the same result every time. This works because the agent takes the model
        as an argument (\code{create_message}) instead of creating it inside; that pattern is
        called \textbf{dependency injection}.
  \item[Environment variables and \code{.env}] Settings a program reads from its environment,
        such as \code{ANTHROPIC_API_KEY} or \code{PM_AGENT_MODEL}. Keeping them out of the code
        means the same code runs with different keys or models, and secrets never get
        committed. \code{.env} is a git-ignored file of \code{KEY=value} lines; a ten-line
        \code{load_dotenv()} copies them into the environment at startup, without overriding
        anything already set in the shell.
  \item[JSON Lines (\code{.jsonl})] A file with one complete JSON object per line. Appending a
        run is just writing one more line, so the file never has to be read and rewritten,
        and a crash can only lose the run in progress. Later code (evals, dashboards) reads it
        line by line.
\end{description}

\subsection{The shell PATH (added in M2b)}

When you type a command such as \code{claude}, the shell (zsh on macOS) looks for a program
with that name in each folder listed in the \textbf{PATH} environment variable, in order.
If none has it, you get ``command not found'', even when the program is installed.

That happened in M2b. The Claude Code installer put \code{claude} in \code{~/.local/bin},
but \code{~/.zshrc} (the script zsh runs at startup) set PATH to a fixed list that did not
include that folder. The fix adds the folder at the front, keeping everything else:

\begin{lstlisting}[style=plain]
export PATH="$HOME/.local/bin:$PATH"
\end{lstlisting}

\code{$PATH} expands to the current value, so this \emph{prepends} rather than replaces.
Writing \code{export PATH="/a:/b"} with no \code{$PATH} is what caused the problem: it throws
away anything added earlier.

\subsection{Git, GitHub, and secrets}

\begin{description}[style=nextline]
  \item[Git] Version control: a history of snapshots (\textbf{commits}) of your files. Each
        commit has a message saying what changed. You can always go back.
  \item[Repository (repo)] A project folder tracked by Git.
  \item[Remote / GitHub] A copy of the repo hosted online. \code{git push} uploads your
        commits; \code{git pull} downloads others'. Ours:
        \url{https://github.com/adityakalkar/portfolio-sentinel}.
  \item[Branch] A separate line of commits. We work on \code{main}.
  \item[.gitignore] A list of files Git should never track: the virtual environment,
        caches, generated results, and above all \code{.env}.
  \item[.env] A file of secret settings (the API key). Never committed. \code{.env.example}
        is the committed template with the values left blank.
\end{description}

\textbf{Public repos.} Anyone can see every committed file \emph{and} the full history. A
secret committed once and deleted later is still visible in history, which is why secrets go
in \code{.env} from day one.

\subsection{CI and GitHub Actions}

\textbf{CI} (continuous integration) runs your tests automatically on every push, on a fresh
machine. It catches ``works on my machine'' problems (a file you forgot to commit, a library
installed locally but missing from \code{pyproject.toml}) and stops broken code from slipping
in unnoticed. A green check or red cross appears next to each commit on GitHub.

\textbf{GitHub Actions} is GitHub's automation service. A \textbf{workflow} is a YAML file in
\code{.github/workflows/}. It has \textbf{triggers} (\code{on:} push, pull request, a
schedule, or manual), \textbf{jobs} that run on a virtual machine (a \textbf{runner}), and
\textbf{steps} inside each job. Steps either \code{uses:} a pre-built action (such as checking
out the code) or \code{run:}s a shell command. It is free for public repos. We use it for CI
now and for the scheduled morning brief in M5.

% =====================================================================
\part{Milestone log}
% =====================================================================

\section{Setup and spec review}

\subsection{What happened}

\begin{enumerate}
  \item Created the repo with the initial project documents: the spec (\code{docs/SPEC.md}),
        the task backlog (\code{docs/TASKS.md}), and \code{CLAUDE.md}.
  \item Read the whole spec as a critical reviewer and raised 8 issues. All were agreed and
        written into \code{docs/SPEC.md}.
  \item Created the GitHub repo \textbf{portfolio-sentinel} with an MIT license,
        description, and topics.
\end{enumerate}

\textbf{What \code{CLAUDE.md} is.} A file Claude Code reads automatically at the start of every
session. It holds the project rules (architecture, conventions, the MCP v2 notes, how to work).
That is why you can \code{/clear} a session without losing the rules.

\textbf{Why MIT.} The MIT license is short and permissive: anyone may read, use, and modify the
code as long as they keep the copyright notice. It is a common choice for open-source projects.

\subsection{The 8 spec decisions and their reasoning}

\begin{enumerate}
  \item \textbf{Leave clear room for the ``\$1M fixes it'' story.} The fix must land clearly
        below the 13.5\% warning line, not on it, so small data changes cannot flip the result.
        This turned out to matter in M1a (Section~\ref{sec:m1a-story}).
  \item \textbf{Flag identity is (rule, subject).} A breach becoming a warning counts as
        ``resolved,'' not ``new.'' Without this rule, \code{simulate_trades} would be ambiguous.
  \item \textbf{``Calm'' means no breaches, not no flags.} The 10/21 brief should still mention
        NVDA's warning in one line.
  \item \textbf{Exact flags per day.} The sanity tests check the exact set of flags, so no other
        limit trips by accident and clutters the story.
  \item \textbf{The made-up-number check exempts some numbers}: numbers from the user's question
        (``\$1M''), dates, and small integers below 10.
  \item \textbf{Manual brief only until real data.} On sample data the latest date never
        changes, so a daily scheduled run would repeat the same brief and spend API credit.
  \item \textbf{Cron at 12:00 UTC on weekdays.} GitHub's schedules run in UTC and ignore daylight
        saving, so 12:00 UTC is 7am (winter) or 8am (summer) Eastern, before market open either
        way.
  \item \textbf{The ``what you could do'' line covers breaches only}, stated as a fact, never as
        advice.
\end{enumerate}

\begin{takeaway}
The spec was reviewed like a design doc before any code was written, and eight issues were
fixed: for example, the planted test scenario sat on a knife edge, and ``new flag'' was
undefined. Fixing ambiguity before coding is cheaper than fixing it in tests or production.
\end{takeaway}

% ---------------------------------------------------------------------
\section{M0: repo setup and CI}

\subsection{Goal}

A project that installs cleanly in a fresh environment, with tests running automatically on
GitHub. Tasks 0.1 (skeleton) and 0.2 (CI).

\subsection{Files created}

\begin{lstlisting}[style=plain]
.github/workflows/ci.yml     the CI workflow
pyproject.toml               project metadata and dependencies
README.md                    placeholder (the real one comes in M7)
src/pm_agent/__init__.py     the package
tests/test_smoke.py          one tiny test
scripts/, data/sample/, evals/, briefs/   folders (with .gitkeep placeholders)
\end{lstlisting}

\textbf{Why \code{.gitkeep}?} Git tracks files, not folders. An empty folder disappears from
the repo. An empty placeholder file called \code{.gitkeep} (a convention, not a Git feature)
keeps it.

\subsection{pyproject.toml, line by line}

\begin{lstlisting}[style=plain]
[build-system]
requires = ["setuptools>=68"]
build-backend = "setuptools.build_meta"
\end{lstlisting}
Which tool builds the package. \textbf{setuptools} is the long-standing standard.

\begin{lstlisting}[style=plain]
[project]
name = "pm-agent"
version = "0.1.0"
readme = "README.md"
license = "MIT"
requires-python = ">=3.11"
dependencies = ["mcp>=2.3,<3", "pandas>=2.2", "numpy>=1.26"]
\end{lstlisting}
The project's identity and needs. \textbf{mcp} is the MCP SDK (for the server in M2).
\textbf{pandas} provides tables (DataFrames) for the analytics. \textbf{numpy} provides fast
numeric arrays and the random number generator. The Anthropic SDK is deliberately \emph{not}
listed yet; it is added in M3, when the agent needs it.

\begin{lstlisting}[style=plain]
[project.optional-dependencies]
dev = ["pytest>=8"]

[tool.setuptools.packages.find]
where = ["src"]

[tool.pytest.ini_options]
pythonpath = ["src"]
testpaths = ["tests"]
\end{lstlisting}
The \code{dev} extra, where to find the package (\code{src/}), and pytest settings: look for
tests in \code{tests/}, and make \code{src/} importable.

\subsection{The CI workflow, line by line}

\begin{lstlisting}[style=plain]
name: CI
on:
  push:
  pull_request:
jobs:
  test:
    runs-on: ubuntu-latest
    steps:
      - uses: actions/checkout@v4          # download the repo onto the runner
      - uses: actions/setup-python@v5      # install Python 3.12
        with:
          python-version: "3.12"
          cache: pip                       # reuse downloaded packages between runs
      - name: Install
        run: pip install -e ".[dev]"
      - name: Run tests
        run: pytest -q
\end{lstlisting}

On every push and pull request, GitHub starts a fresh Ubuntu Linux machine, checks out the
code, installs Python and the project, and runs the tests. If any step fails, the run is red.
The \code{@v4} and \code{@v5} pin the version of each pre-built action, so an update to the
action cannot silently change our build.

\subsection{Decisions}

\begin{itemize}
  \item \textbf{One smoke test instead of zero.} The task said ``0 tests is fine,'' but pytest
        exits with code 5 when it finds no tests, and CI treats any non-zero exit code as a
        failure. A one-line test that imports the package keeps CI green and proves the
        install works.
  \item \textbf{A placeholder README.} \code{pyproject.toml} points to \code{README.md}, and the
        install fails if it is missing. It also gives the GitHub page some content.
  \item \textbf{pandas 3.1 was installed.} pandas 3 enables ``copy-on-write'': a DataFrame
        derived from another no longer silently changes the original when modified. This helps
        with our rule that what-if trades must not change the original portfolio. Examples
        online written for pandas 2 may behave slightly differently.
\end{itemize}

\subsection{How it was verified}

\begin{enumerate}
  \item Created a fresh virtual environment, ran \code{pip install -e ".[dev]"}, and ran
        pytest: \code{1 passed}.
  \item Checked that \code{ci.yml} is valid YAML.
  \item Pushed. The first GitHub Actions run succeeded on every step (checkout, set up Python,
        install, run tests).
\end{enumerate}

\begin{reviewq}
\begin{enumerate}
  \item Why does \code{.venv/} go in \code{.gitignore}, while \code{pyproject.toml} is committed?
  \item What does \code{mcp>=2.3,<3} allow, and why exclude version 3?
  \item What would happen in CI if we had zero tests?
  \item What does the \code{-e} in \code{pip install -e} do?
  \item Name two kinds of mistake CI catches that running tests locally can miss.
\end{enumerate}
\end{reviewq}

% ---------------------------------------------------------------------
\section{M1a: sample data and loaders}

\subsection{Goal}

Tasks 1.1 and 1.2: generate a realistic but synthetic fund whose data contains a known story,
and write \code{data.py}, the only module that reads data files.

\subsection{Why synthetic data with a planted story?}

To test whether the morning brief ``finds what matters,'' you need days where you
\emph{know} what matters. Real market data has no answer key. So we generate data where we
planted the events ourselves, and we can check whether the agent catches them. This is a common
technique in testing AI systems: control the inputs so the expected output is known.

\subsection{The four data files}

\begin{lstlisting}[style=plain]
data/sample/
  portfolios/sample_fund.json        display_name, benchmark, cash, limits, positions
  benchmarks/sample_large_cap.csv    ticker, weight      (23 stocks, sums to 1)
  prices.csv                         date, ticker, close (25 tickers x 5 dates)
  securities.csv                     ticker, name, sector
\end{lstlisting}

\textbf{CSV} (comma-separated values) is a plain-text table. \textbf{JSON} is a plain-text
format for nested data. We use JSON for the portfolio because it has nested parts (a limits
dictionary and a list of positions).

\subsection{The fund and its benchmark}

The \textbf{Sample Large Cap Fund}: \$50M NAV on 2026-10-21, 22 positions across 8 sectors,
4\% cash, and one short (INTC, $-2\%$). The benchmark is a 23-stock stand-in for the S\&P~500.
On purpose, it excludes two of the fund's holdings (AMD and INTC) and includes three stocks the
fund does not hold (ORCL, ABBV, KO), so the code must handle both ``held but not in the
benchmark'' and ``in the benchmark but not held.''

Sector comparison on 2026-10-21 (fund net weight versus benchmark):

\begin{center}
\begin{tabular}{@{}lrrr@{}}
\toprule
\textbf{Sector} & \textbf{Fund} & \textbf{Benchmark} & \textbf{Active} \\
\midrule
Information Technology & 37.0\% & 40.0\% & $-3.0\%$ \\
Communication Services & 14.0\% & 14.0\% & 0.0\% \\
Consumer Discretionary & 10.5\% & 12.0\% & $-1.5\%$ \\
Health Care            & 9.5\%  & 10.0\% & $-0.5\%$ \\
Financials             & 11.5\% & 11.0\% & $+0.5\%$ \\
Consumer Staples       & 5.0\%  & 5.0\%  & 0.0\% \\
Industrials            & 4.5\%  & 4.0\%  & $+0.5\%$ \\
Energy                 & 4.0\%  & 4.0\%  & 0.0\% \\
\midrule
Total                  & 96.0\% & 100.0\% & $-4.0\%$ \\
\bottomrule
\end{tabular}
\end{center}

The fund total is 96\% because 4\% is cash. Tech is 37\% net: 39\% long (NVDA, AAPL, MSFT,
AVGO, AMD) minus the 2\% INTC short.

\subsection{The planted story}\label{sec:m1a-story}

\begin{center}
\begin{tabular}{@{}lll@{}}
\toprule
\textbf{Date} & \textbf{NVDA weight} & \textbf{Status vs 15\% limit} \\
\midrule
10/16--10/20 & 13.94\%--14.02\% & warning (at or above 13.5\%) \\
10/21        & 14.00\% & warning \\
10/22        & \textbf{15.20\%} & \textbf{breach} \\
10/22 after selling \$1M & 13.23\% & clear \\
\bottomrule
\end{tabular}
\end{center}

On 2026-10-22: NVDA $+10\%$, UNH $-8\%$ (the biggest loser), AMD $+4\%$, INTC $-3\%$ (the short
gains). Every other stock moves less than 0.5\%, so nothing competes with the story.

\textbf{Why NVDA starts at 14.0\%, not the 14.2\% in the original spec.} The arithmetic:

\begin{itemize}
  \item NVDA's value on 10/22 $= 7{,}000{,}020 \times 1.10 = \$7{,}700{,}022$.
  \item The fund's NAV rises to about \$50.65M (NVDA's gain, minus UNH's loss, plus AMD and the
        INTC short).
  \item NVDA weight $= 7.70 / 50.65 \approx 15.20\%$: a breach.
  \item After a \$1M sale, NAV is unchanged (stock becomes cash), and NVDA's weight is
        $6.70 / 50.65 \approx 13.23\%$, clearly under 13.5\%.
\end{itemize}

At 14.2\%, the same sale would have landed at about 13.44\%, almost on the line. That is the
``knife edge'' from the spec review, decision~1. Moving the starting weight to 14.0\% keeps a clear
breach \emph{and} a clear fix.

\textbf{Making sure nothing else trips.} The benchmark weights were chosen so that, on every
day, the largest stock active weight is about 2.5--3.2\% (warning line 4.5\%), the largest
sector active weight is about 3\% (warning line 7.2\%), gross is 100\% (warning line 103.5\%),
and net is 96\% (warnings below 92\% or above 100\%). So the only flag is NVDA.

\subsection{How the generator works}

\code{scripts/make_sample_data.py} keeps every story number in clearly named constants at the
top of the file (\code{FUND_TARGET_WEIGHTS}, \code{BENCHMARK_WEIGHTS}, \code{STORY_MOVES},
\code{LIMITS}), so the story is visible just by reading them.

\textbf{Step 1: prices.} Each stock has an \emph{anchor price} on 10/21. Earlier days are a
small random walk \emph{backwards} from the anchor:

\begin{lstlisting}[style=py]
price = anchor
for d in reversed(DATES[:anchor_i]):          # 10/20, 10/19, 10/16
    price = price / (1 + rng.normal(0, DAILY_NOISE))
    prices[d][ticker] = round(price, 2)
\end{lstlisting}

Going backwards keeps 10/21 exactly at the anchor, so the target weights are exact on that day.
On 10/22, the planted stocks get their story move and everyone else gets tiny noise.

\textbf{Step 2: shares.} Shares $=$ target weight $\times$ NAV $/$ anchor price, rounded to whole
shares (as in a real fund). Cash is whatever makes NAV exactly \$50M, which turned out to be
\$2{,}000{,}074, about 4\%. Weights therefore match their targets to within 0.01\%.

\textbf{Determinism and the random seed.} A computer's ``random'' numbers come from a formula
started by a \textbf{seed}. The same seed always gives the same sequence. We fix
\code{SEED = 42}, so the generator produces identical files every run. Tests depend on exact
numbers, so they must not change from run to run.

A subtle detail: the noise for 10/22 is drawn for \emph{every} stock, even the ones whose story
move replaces it:

\begin{lstlisting}[style=py]
noise = rng.normal(0, STORY_DAY_NOISE)  # always drawn, so the random sequence is stable
move = STORY_MOVES.get(ticker, noise)
\end{lstlisting}

If we only drew noise for non-story stocks, adding or removing a story stock would shift every
random number after it and change unrelated prices.

\subsection{data.py: the loaders}

\begin{center}
\begin{tabularx}{\textwidth}{@{}l X@{}}
\toprule
\textbf{Function} & \textbf{What it returns} \\
\midrule
\code{data_dir()}            & the data folder: \code{PM_AGENT_DATA_DIR} if set, else \code{data/sample} \\
\code{list_portfolios()}     & portfolio names, e.g. \code{["sample_fund"]} \\
\code{load_portfolio(name)}  & dict: name, display\_name, benchmark, cash, limits, positions (DataFrame) \\
\code{load_benchmark(name)}  & DataFrame of ticker, weight \\
\code{load_securities()}     & DataFrame of ticker, name, sector \\
\code{load_prices()}         & DataFrame of date, ticker, close \\
\code{available_dates()}     & sorted list of price dates \\
\code{prices_on(date=None)}  & (date used, prices by ticker); defaults to the latest date \\
\code{previous_date(date)}   & the price date before \code{date} (skips weekends) \\
\bottomrule
\end{tabularx}
\end{center}

Design decisions worth explaining:

\begin{itemize}
  \item \textbf{Helpful errors.} An unknown name or date raises \code{ValueError} with a list of
        what \emph{is} available, for example ``No prices for 2026-10-17. Available dates:
        2026-10-16, \ldots''. This is aimed at the AI: in M2 these messages are passed to the
        model, which can read them and retry with a valid date. Good error messages are how
        agents correct themselves.
  \item \textbf{The env var is read on every call}, not once at import. Tests and evals change
        \code{PM_AGENT_DATA_DIR} while running, and a value captured at import time would ignore
        that.
  \item \textbf{Dates stay as strings} (\code{"2026-10-22"}). pandas would otherwise convert them
        to timestamps. \code{YYYY-MM-DD} strings sort correctly as text, compare simply, and are
        what the AI sends and receives.
  \item \textbf{\code{load_portfolio} returns a plain dict for now.} The \code{Portfolio} class
        belongs to the analytics module (task 1.3), so it wasn't built early. In M1b this
        function will return a \code{Portfolio}.
  \item \textbf{A pandas Series for prices.} A Series is a one-column table with labels (here,
        tickers), so \code{prices["NVDA"]} looks up a price directly.
\end{itemize}

\subsection{The tests}

\code{tests/test_data.py} has 13 tests. Three examples:

\begin{lstlisting}[style=py]
def test_unknown_date_lists_available():
    with pytest.raises(ValueError, match="No prices for 2026-10-17.*2026-10-16.*2026-10-22"):
        data.prices_on("2026-10-17")
\end{lstlisting}
Checks that the error happens \emph{and} that its message lists the available dates.
2026-10-17 is a Saturday, so it has no prices.

\begin{lstlisting}[style=py]
def test_previous_date():
    assert data.previous_date("2026-10-22") == "2026-10-21"
    assert data.previous_date("2026-10-19") == "2026-10-16"  # skips the weekend
\end{lstlisting}
``Previous date'' means the previous \emph{trading} day, not the previous calendar day.

\begin{lstlisting}[style=py]
def test_data_dir_override(tmp_path, monkeypatch):
    ...write a tiny portfolio and prices file into tmp_path...
    monkeypatch.setenv("PM_AGENT_DATA_DIR", str(tmp_path))
    assert data.list_portfolios() == ["tiny"]
\end{lstlisting}
Uses two pytest fixtures: \code{tmp_path} (a throwaway folder) and \code{monkeypatch} (sets the
environment variable for this test only). This proves the override works, which evals will rely
on later.

\subsection{How it was verified}

\begin{enumerate}
  \item Ran the generator twice and compared the outputs with \code{diff}: identical.
  \item Checked the story on the generated files with a throwaway script (the table in
        Section~\ref{sec:m1a-story}). The permanent version of this check is task 1.8 in M1b.
  \item Watched the new tests fail before \code{data.py} existed, then pass after: 14 tests
        passing in total.
\end{enumerate}

\begin{reviewq}
\begin{enumerate}
  \item Why plant a story in synthetic data instead of using real prices?
  \item NVDA is 14\% of a \$50M fund and rises 10\% while the fund's NAV rises about 1.3\%. Why
        does NVDA's weight rise to about 15.2\%, not exactly 15.4\%?
  \item Why does selling \$1M of NVDA leave NAV unchanged?
  \item What is a random seed, and why does the generator fix one?
  \item Why does the generator walk prices \emph{backwards} from 10/21?
  \item Why do \code{data.py}'s error messages list what is available? Who is the real audience
        for those messages?
  \item What would break if \code{data_dir()} read the environment variable only once, at import?
\end{enumerate}
\end{reviewq}

\begin{takeaway}
The synthetic data has a planted scenario so the agent's output can be scored against a
known answer. Building it revealed that the spec's own numbers put the key test case within
0.06 percentage points of a threshold, so the starting weight was adjusted to make the test
robust, and the change was documented.
\end{takeaway}

% ---------------------------------------------------------------------
\section{M1b: the analytics core}

\subsection{Goal}

Tasks 1.3 to 1.8: write \code{src/pm_agent/analytics.py}, the one place where portfolio math
happens, test-first; then add sanity tests proving the sample data still tells its planted
story. Everything the agent will ever say about numbers traces back to this file, so it has to
be right and provably so.

\subsection{What was built}

\begin{center}
\begin{tabularx}{\textwidth}{@{}l l X@{}}
\toprule
\textbf{Task} & \textbf{Function} & \textbf{What it returns} \\
\midrule
1.3 & \code{Portfolio} & dataclass: name, positions (ticker $\to$ shares), cash, benchmark, limits \\
1.3 & \code{position_weights} & table: ticker, shares, price, market value, weight; biggest first \\
1.3 & \code{exposure_summary} & NAV, cash, long/short/gross/net exposure, position counts \\
1.4 & \code{active_weights} & per stock in fund \emph{or} benchmark: portfolio, benchmark, active weight \\
1.4 & \code{sector_exposure} & per sector: long, short, net, gross (+ benchmark, active) \\
1.5 & \code{check_limits} & list of flags, breaches first; \code{[]} means all clear \\
1.6 & \code{compare_flags} & which flags a change resolved and which it created \\
1.6 & \code{apply_trades} & a \emph{new} Portfolio after dollar trades \\
1.7 & \code{weight_changes} & per position: start/end weight, change, price return \\
1.8 & (tests only) & \code{tests/test_sample_data.py}: the planted story holds \\
\bottomrule
\end{tabularx}
\end{center}

Also: \code{data.load_portfolio} now returns a \code{Portfolio} instead of a plain dict, as
promised in M1a.

\textbf{Test counts.} 57 tests in total: 34 known-answer tests in \code{tests/test_analytics.py},
9 sanity tests in \code{tests/test_sample_data.py}, and the 14 from earlier milestones.

\subsection{The shape of every function}

Every analytics function takes its inputs as arguments and returns a value. Nothing is loaded
inside. The inputs are always the same few kinds of thing:

\begin{lstlisting}[style=py]
portfolio: Portfolio      # what we hold
prices:    pd.Series      # ticker -> close on one date
sectors:   pd.Series      # ticker -> sector name
benchmark: pd.Series      # ticker -> benchmark weight
\end{lstlisting}

This is what ``pure'' buys us. A test can build a three-stock portfolio with \$1 prices in two
lines and check the answer by hand, with no files at all:

\begin{lstlisting}[style=py]
port = Portfolio(name="t", positions={"LONG": 10, "SHORT": -2}, cash=20)
s = exposure_summary(port, prices(LONG=10, SHORT=10))
assert s["gross_exposure"] == pytest.approx(1.20)   # $100 long, $20 short, NAV $100
assert s["net_exposure"] == pytest.approx(0.80)
\end{lstlisting}

The M2 server will do the wiring: load with \code{data.py}, call analytics, convert to percent.

\textbf{Which module depends on which.} \code{Portfolio} is defined in \code{analytics.py}, and
\code{data.py} imports it to build one from JSON. Analytics imports nothing from the project,
so it stays pure and can be tested on its own. The arrow only points one way:
\code{data} $\to$ \code{analytics}.

\subsection{Weights and exposure (1.3)}

\begin{lstlisting}[style=py]
def position_weights(portfolio, prices):
    missing = sorted(t for t in portfolio.positions if t not in prices.index)
    if missing:
        raise ValueError(f"No price for {', '.join(missing)} on this date")
    rows = pd.DataFrame({"ticker": list(portfolio.positions),
                         "shares": list(portfolio.positions.values())})
    rows["price"] = rows["ticker"].map(prices).astype(float)
    rows["market_value"] = rows["shares"] * rows["price"]
    nav = rows["market_value"].sum() + portfolio.cash
    rows["weight"] = rows["market_value"] / nav
    ...sort by absolute market value, biggest first...
\end{lstlisting}

Line by line: check every holding has a price (fail loudly, naming the missing tickers,
rather than silently treating a missing price as \$0); look up each price with \code{.map};
market value is shares times price; NAV is the sum plus cash; weight is value over NAV.
Sorting by \emph{absolute} value puts a big short near the top, where a reader expects it.

\code{exposure_summary} builds on it: long is the sum of positive weights, short the sum of
negative ones, then gross $=$ long $-$ short and net $=$ long $+$ short.

On the sample fund on 2026-10-22 it gives NAV \$50{,}646{,}833, long 97.97\%, short
$-1.92\%$, gross 99.88\%, net 96.05\%, with 21 longs and 1 short.

\subsection{Benchmark comparison (1.4)}

\code{active_weights} relies on pandas \emph{index alignment} (Part~I):

\begin{lstlisting}[style=py]
held = position_weights(portfolio, prices).set_index("ticker")["weight"]
table = pd.DataFrame({"portfolio_weight": held, "benchmark_weight": benchmark})
table = table.fillna(0.0)
table["active_weight"] = table["portfolio_weight"] - table["benchmark_weight"]
\end{lstlisting}

Putting two Series side by side gives one row per ticker in either. A stock only the fund holds
(AMD) gets a benchmark weight of 0, so its whole weight is active. A stock only the benchmark
holds (ORCL) gets a portfolio weight of 0, so it is underweight by its full benchmark weight.
The result is sorted by absolute active weight: biggest bets first, in either direction.

\code{sector_exposure} splits each weight into its long and short part with \code{clip}, adds
them up by sector with \code{groupby}, and derives net and gross. With a benchmark, it also adds
up the benchmark's weights by sector and computes active $=$ net $-$ benchmark. A sector that
only the benchmark holds still gets a row (all zeros for the fund), because being absent from a
sector is a sector bet.

Known-answer test, a sector with a short: Tech holds A at $+50\%$ and C short at $-10\%$, so
Tech long $= 50\%$, short $= -10\%$, net $= 40\%$, gross $= 60\%$.

\subsection{Limit checks (1.5)}

\code{check_limits} works in three steps.

\textbf{1. Validate the limits.} Unknown keys raise an error listing the known ones. A typo
such as \code{max_positon} would otherwise be silently ignored, and a limit you think is
being checked would not be. The active-weight rules also require a benchmark (and the sector
rule requires sectors); if one is missing, that is an error too, for the same reason.

\textbf{2. Collect every \code{max_*} check, then rate each one} with one small helper:

\begin{lstlisting}[style=py]
def _max_rule_status(value, limit, warn_at):
    size = abs(value)
    if size > limit:
        return "breach"
    if size >= warn_at * limit:
        return "warning"
    return None
\end{lstlisting}

Position weights, stock active weights, sector active weights and gross exposure all go
through this same function, so the warning rule is written exactly once.

\textbf{3. Net rules use the fixed buffer instead.} \code{min_net} breaches below 90\% and warns
up to 92\%; \code{max_net} breaches above 102\% and warns from 100\%.

Finally, flags are sorted with breaches first. Because Python's \code{sorted} is stable, the
order within each group is the order the checks ran in (positions, then stock active, sector
active, gross, net), and within a rule the biggest first.

\textbf{Worked example} (from the tests). NAV \$100 with A $= 16\%$, B $= -14\%$ (short),
C $= 10\%$; position limit 15\%, warning line 13.5\%:

\begin{center}
\begin{tabular}{@{}lrll@{}}
\toprule
\textbf{Ticker} & \textbf{Weight} & \textbf{Compared size} & \textbf{Status} \\
\midrule
A & $16\%$  & $16\% > 15\%$ & breach \\
B & $-14\%$ & $14\% \ge 13.5\%$ & warning (shorts count by size) \\
C & $10\%$  & $10\% < 13.5\%$ & no flag \\
\bottomrule
\end{tabular}
\end{center}

\textbf{The net buffer, tested at every edge}, with \code{warn_at} $= 0.9$ deliberately also set
to prove net ignores it:

\begin{center}
\begin{tabular}{@{}rl@{}}
\toprule
\textbf{Net} & \textbf{Result} \\
\midrule
89\%  & \code{min_net} breach \\
91\%  & \code{min_net} warning \\
95\%  & no flag (with \code{warn_at}, it would be a \code{max_net} warning: $95 \ge 91.8$) \\
101\% & \code{max_net} warning \\
103\% & \code{max_net} breach \\
\bottomrule
\end{tabular}
\end{center}

\subsection{What-if trades and flag comparison (1.6)}

\begin{lstlisting}[style=py]
def apply_trades(portfolio, prices, trades):
    positions = dict(portfolio.positions)  # a copy, so the original is never modified
    cash = portfolio.cash
    for trade in trades:
        ticker, dollars = trade["ticker"], trade["dollars"]
        if ticker not in prices.index:
            raise ValueError(f"Unknown ticker '{ticker}'. Tickers with prices: ...")
        shares = float(positions.get(ticker, 0) + dollars / prices[ticker])
        if math.isclose(shares, 0, abs_tol=1e-9):
            positions.pop(ticker, None)
        else:
            positions[ticker] = shares
        cash -= dollars
    return replace(portfolio, positions=positions, cash=cash)
\end{lstlisting}

The whole function is ``shares change by dollars over price, cash changes by minus dollars''.
\code{positions.get(ticker, 0)} means a ticker not yet held starts at 0 shares, so the same
line both adds to a position and opens a new one. Selling more than you hold leaves negative
shares: a short, consistent with the conventions. The copy-then-\code{replace} pattern is what
guarantees the original portfolio is untouched (Part~I, dataclasses), and a test checks it.

\code{compare_flags(before, after)} implements the spec's flag identity rules. It indexes both
lists by \code{(rule, subject)} and then:

\begin{center}
\begin{tabular}{@{}lll@{}}
\toprule
\textbf{Before} & \textbf{After} & \textbf{Counted as} \\
\midrule
breach  & warning  & resolved \\
any     & (no flag) & resolved \\
warning & breach   & new \\
(no flag) & any    & new \\
same status & same status & neither \\
\bottomrule
\end{tabular}
\end{center}

The headline test, ``trimming fixes a breach'': A is 16\% of a \$100 fund; selling \$2 brings it
to 14\%. That is still a warning, but \code{compare_flags} reports the breach as resolved and
nothing as new.

\subsection{Weight changes (1.7)}

\code{weight_changes} holds the shares fixed, computes weights on two dates, and subtracts. It
also returns each stock's price return. The known-answer test: two \$50 positions, no cash; A's
price doubles, so A is now \$100 of \$150. Weights go from 50/50 to 66.7/33.3, although nobody
traded. This is drift (Part~I).

\subsection{Sample data sanity tests (1.8)}

\code{tests/test_sample_data.py} checks the real sample files, end to end through
\code{data.py} and analytics:

\begin{itemize}
  \item The benchmark sums to 1. Every held or benchmark ticker has a sector and a price on
        every date.
  \item On each of 10/16, 10/19, 10/20 and 10/21, the \emph{exact} flag list is one NVDA
        position warning. On 10/22 it is exactly one NVDA breach (15.20\%). No other limit
        trips anywhere (spec decision 4).
  \item 10/22 price moves: NVDA $+10\%$, UNH $-8\%$ (the biggest loser), AMD $+4\%$, INTC
        $-3\%$; NVDA has the biggest weight drift.
  \item Selling \$1M of NVDA on 10/22 brings it to 13.23\%, at least 0.2 points under the
        13.5\% warning line; afterwards there are no flags at all, and \code{compare_flags}
        reports exactly one resolved breach and nothing new.
\end{itemize}

An \code{autouse} fixture removes \code{PM_AGENT_DATA_DIR} for these tests, so they always
check the checked-in sample data even if the variable is set in the shell.

\subsection{Decisions and why}

\begin{enumerate}
  \item \textbf{Analytics takes data as arguments (Series), not names or paths.} This keeps it
        pure and makes known-answer tests tiny. The alternative, having analytics call
        \code{data.py}, would mean every test needs data files.
  \item \textbf{Positions are a plain \code{dict} (ticker $\to$ shares), not a DataFrame.}
        Easier to read, to build in tests, and to change in \code{apply_trades}. Tables are
        built from it when needed.
  \item \textbf{Flags carry the signed value.} The check uses the size, but the reader needs
        the direction ($-14\%$ short versus $+14\%$ long; underweight versus overweight).
  \item \textbf{Exactly at the limit is a warning.} ``Max 15\%'' allows 15\%. The test for D
        (active exactly $-5\%$) pins this down.
  \item \textbf{Missing settings disable the warning band.} If \code{warn_at} is absent it
        defaults to 1.0 (warn only at exactly the limit); if \code{net_warn_buffer} is absent,
        0. No made-up thresholds appear that the portfolio file did not ask for.
  \item \textbf{A configured active limit with no benchmark is an error}, not silently
        skipped. Silently skipping a configured limit is the most dangerous possible
        failure for a risk check.
  \item \textbf{\code{compare_flags} lives in analytics.} The spec puts resolved/new flags in
        the \code{simulate_trades} tool output, but the comparison is logic, and the MCP layer
        must not contain logic (architecture rule 4). It was not named in the task list, so it
        was added in 1.6, where its first test (``trimming fixes a breach'') lives.
  \item \textbf{Sort orders.} Positions by absolute value, active weights and weight changes by
        absolute change, sectors by gross exposure (largest first). ``Biggest first'' is what
        a PM scans for, and it lets \code{top_n} in M2 simply take the first rows.
  \item \textbf{Fractional shares from dollar trades}, with a float-tolerance check so a fully
        sold position is removed rather than left at $10^{-13}$ shares.
  \item \textbf{Shares are stored as plain \code{float}}, not numpy numbers, so they convert
        cleanly to JSON at the MCP boundary later.
\end{enumerate}

\subsection{How it was verified}

\begin{enumerate}
  \item For each of 1.3 to 1.7: wrote the tests, ran them and watched them fail (the function
        did not exist yet), implemented it, and watched them pass.
  \item For 1.8, the sanity tests passed on first run because the data was already right, so
        a mutation check was done (Part~I): NVDA's 10/22 price was changed from \$198.00 to
        \$190.00, three tests failed, and the file was restored with \code{git checkout}.
  \item Ran the analytics on the real sample data before writing the sanity tests, to see the
        actual flags on every date (they matched the M1a story table exactly).
  \item Full suite: 57 passed.
\end{enumerate}

\begin{reviewq}
\begin{enumerate}
  \item Why does \code{check_limits} compare $|\text{value}|$ with the limit, but report the
        signed value?
  \item A position is at exactly 15.0\% with a 15\% limit. Breach, warning, or neither? Why?
  \item Why do net limits use a fixed 2-point buffer instead of \code{warn_at}? Give the net
        exposure that shows the difference.
  \item What does \code{.fillna(0.0)} do in \code{active_weights}, and which two situations does
        it handle?
  \item \code{Portfolio} is frozen. Why must \code{apply_trades} still copy the positions dict?
  \item A trade turns an NVDA breach into an NVDA warning. Is that a resolved flag, a new flag,
        or neither? What if a warning turns into a breach?
  \item Why is \code{compare_flags} in \code{analytics.py} and not in the MCP server?
  \item Why might a sanity test that passes on its first run still be worthless, and how do you
        check it is not?
  \item What could go wrong if \code{check_limits} silently skipped unknown limit keys?
\end{enumerate}
\end{reviewq}

\begin{takeaway}
All the portfolio math now lives in one pure, tested module. Every rule from the spec has a
known-answer test small enough to check by hand, including the edges (exactly at a limit, shorts,
stocks missing from one side, the net buffer). Wherever a mistake could hide a risk, such as a
misspelled limit or a missing benchmark, the code raises an error instead of guessing.
Sanity tests then pin the planted story to the real sample data, and a deliberate mutation
showed that those tests catch a broken story.
\end{takeaway}

% ---------------------------------------------------------------------
\section{M2a: the MCP server}

\subsection{Goal}

Task 2.1: wrap the analytics as 8 MCP tools (SPEC~\S7) so that any MCP client (Claude Code
in task 2.2, our own agent in M3) can ask portfolio questions and get exact numbers back.
Every tool is read-only, returns structured output, and turns bad input into an error message
the model can read and act on.

\subsection{What was built}

\code{src/pm_agent/mcp_server.py}, plus two small loaders in \code{data.py}
(\code{benchmark_weights(name)} and \code{sector_map()}, which return the Series shape the
analytics functions take).

\begin{center}
\begin{tabularx}{\textwidth}{@{}l l X@{}}
\toprule
\textbf{Tool} & \textbf{Calls (analytics)} & \textbf{Example question} \\
\midrule
\code{list_portfolios}        & (data only)         & Which funds can I look at? \\
\code{get_portfolio_summary}  & \code{exposure_summary} & What's our NAV and net exposure? \\
\code{get_positions}          & \code{position_weights} & What are our five biggest positions? \\
\code{compare_to_benchmark}   & \code{active_weights}   & Where are we most overweight? \\
\code{get_sector_exposure}    & \code{sector_exposure}  & What are our sector bets? \\
\code{check_limits}           & \code{check_limits}     & Are we breaching any limits? \\
\code{simulate_trades}        & \code{apply_trades}, \code{check_limits}, \code{compare_flags} & If I sell \$1M of NVDA, do we still breach anything? \\
\code{get_weight_changes}     & \code{weight_changes}   & What moved overnight? \\
\bottomrule
\end{tabularx}
\end{center}

\subsection{How one tool is put together}

Every tool has the same three-step shape: \textbf{load} (data.py), \textbf{compute}
(analytics.py), \textbf{format} (percent and rounding). Here is \code{check_limits}:

\begin{lstlisting}[style=py]
@tool
def check_limits(portfolio: str = DEFAULT_PORTFOLIO, date: str | None = None) -> dict[str, Any]:
    """Check the fund against its risk limits. ...
    Example: "Are we breaching any limits?" or "Is anything close to a limit?"
    """
    with readable_errors():
        port, used_date, prices = load(portfolio, date)                     # 1. load
        flags = analytics.check_limits(port, prices, data.sector_map(),
                                       benchmark_for(port))                  # 2. compute
        return {                                                             # 3. format
            "portfolio": port.name,
            "date": used_date,
            "breach_count": sum(f["status"] == "breach" for f in flags),
            "warning_count": sum(f["status"] == "warning" for f in flags),
            "limits": limit_fields(port.limits),
            "flags": [flag_fields(f) for f in flags],
        }
\end{lstlisting}

The supporting pieces, each a few lines:

\begin{description}[style=nextline]
  \item[\code{@tool}] Registers the function with the server as read-only. The docstring
        becomes the description (see ``What we learned about the SDK'' below).
  \item[\code{readable_errors()}] A context manager: any \code{ValueError} inside the
        \code{with} block becomes a \code{ToolError}, so the message reaches the model.
  \item[\code{pct(weight)}] \code{round(weight * 100, 2)}: 0.152034 becomes 15.2. This is
        a unit conversion, not new math, which is what the architecture rules allow here.
  \item[\code{load(portfolio, date)}] Returns the portfolio, the date actually used (the
        latest one if none was given), and that day's prices. Every tool reports the date it
        used, so the model never has to guess which day a number is from.
\end{description}

\textbf{What the model receives} for ``Are we breaching any limits?'' on 2026-10-22:

\begin{lstlisting}[style=plain]
{
  "portfolio": "sample_fund", "date": "2026-10-22",
  "breach_count": 1, "warning_count": 0,
  "limits": {"max_position_pct": 15.0, "max_stock_active_pct": 5.0,
             "max_sector_active_pct": 8.0, "max_gross_pct": 115.0,
             "min_net_pct": 90.0, "max_net_pct": 102.0,
             "warn_at": 0.9, "net_warn_buffer_pct": 2.0},
  "flags": [{"rule": "max_position", "subject": "NVDA",
             "value_pct": 15.2, "limit_pct": 15.0, "status": "breach"}]
}
\end{lstlisting}

And for a mistake, \code{portfolio="nope"}:

\begin{lstlisting}[style=plain]
Error executing tool check_limits: Unknown portfolio 'nope'. Available: sample_fund
\end{lstlisting}

The model can read ``Available: sample\_fund'' and retry with the right name. That message was
written in M1a (\code{data.py}); the MCP layer just lets it through.

\subsection{The what-if tool in detail}

\code{simulate_trades} is the richest tool. For ``If I sell \$1M of NVDA, do we still breach
anything?'', it returns the before and after summary and flags, plus:

\begin{lstlisting}[style=plain]
"resolved_flags": [{"rule": "max_position", "subject": "NVDA", "value_pct": 15.2, ...}],
"new_flags": [],
"traded_positions": [{"ticker": "NVDA", "shares_before": 38889.0,
                      "shares_after": 33838.49,
                      "weight_before_pct": 15.2, "weight_after_pct": 13.23}]
\end{lstlisting}

\textbf{Why \code{traded_positions} exists.} After the sale, NVDA has no flag at all, so the
``after'' flags cannot tell the model NVDA's new weight. But the morning brief (SPEC~\S9) has
to say ``selling \$1M of NVDA would bring it to 13.23\%'', and every number must come from a
tool. So the tool returns each traded stock's weight before and after. Without it, the model
would be tempted to work out 13.23\% itself, which is exactly what the project forbids.

\subsection{Units, rounding, and defaults}

\begin{itemize}
  \item \textbf{Percent fields} (\code{*_pct}) are rounded to 2 decimals. Differences between
        percentages (\code{active_pct}, \code{weight_change_pct}) are percentage points;
        the server instructions say so, so the model does not call a 1.2-point change
        ``1.2\% growth''.
  \item \textbf{Dollar fields} (\code{*_usd}) are rounded to cents.
  \item \textbf{\code{warn_at} stays a fraction} (0.9), because it is not a weight: it scales
        each limit (0.9 $\times$ 15\% $=$ 13.5\%). Every other limit is converted to percent.
  \item \textbf{Defaults:} portfolio \code{sample_fund}; date the latest price date;
        \code{top_n} 10 (0 means all); weight changes default to overnight (the price date
        before the end date, to the end date).
  \item \textbf{List tools report the full count} (\code{total_positions},
        \code{total_stocks}), so with the default \code{top_n} of 10 the model knows it is
        seeing 10 of 22, not everything.
\end{itemize}

\subsection{What we learned about the SDK}

Before writing the server, a throwaway script checked the installed SDK (\code{mcp} 2.3.0)
against the notes in \code{CLAUDE.md}. Everything matched (\code{MCPServer},
\code{dict[str, Any]} for structured output, \code{ToolError}, the in-process \code{Client},
snake\_case fields). Three details were new:

\begin{enumerate}
  \item \textbf{A \code{ToolError}'s message arrives prefixed}: ``Error executing tool
        check\_limits: Unknown portfolio \ldots''. A plain \code{ValueError} arrives as just
        ``Error executing tool check\_limits'', with the useful part gone.
  \item \textbf{Docstrings are sent with their code indentation}, so every line after the
        first started with four spaces. The \code{@tool} helper runs \code{inspect.cleandoc}
        on the docstring first, which removes it.
  \item \textbf{\code{Args:} sections in docstrings are not turned into parameter
        descriptions.} So each description explains its parameters in plain sentences, and
        the trades parameter uses a \code{TypedDict} so the JSON Schema names its fields
        (\code{ticker}, \code{dollars}). A bare \code{list[dict]} would have produced a
        schema that says only ``array of objects''.
\end{enumerate}

\subsection{Decisions and why}

\begin{enumerate}
  \item \textbf{A \code{with readable_errors():} block in each tool, rather than a wrapping
        decorator.} A decorator that wraps the function can hide its signature from the SDK,
        which builds the input schema from it. The \code{with} block is explicit and safe.
  \item \textbf{\code{TypedDict} instead of a pydantic model for trades.} pydantic is
        installed (the MCP SDK depends on it), but importing it directly would make it our
        dependency without listing it. \code{TypedDict} is in the standard library.
  \item \textbf{Default \code{top_n} of 10, not ``all''.} With real data in M6 the benchmark
        has about 500 stocks; dumping all of them costs tokens and buries the answer. The
        total count keeps the model honest about what it has not seen.
  \item \textbf{\code{simulate_trades} returns traded positions' weights}, so the
        ``would bring it to 13.23\%'' sentence in the brief comes from a tool, not the model.
  \item \textbf{Counting flags in the server} (\code{breach_count}) is tallying, not finance
        math, so it stays in the formatting layer. It saves the model from counting a list.
  \item \textbf{Two new loaders in \code{data.py}}, because both the server and the sanity
        tests needed the same reshaping (benchmark and sectors as Series). Keeping it in
        \code{data.py} means one place to change when real data arrives in M6.
  \item \textbf{Names and sectors included in rows}, so ``what's our biggest health care
        position?'' needs one call, not two.
\end{enumerate}

\subsection{How it was verified}

\begin{enumerate}
  \item Wrote \code{tests/test_mcp_server.py} (16 tests) first and watched it fail on the
        missing module; then implemented the server until all passed. They cover the four
        ``done when'' checks from TASKS.md: all 8 tools listed with descriptions (and
        read-only), \code{check_limits} returning the NVDA breach in percent, the \$1M sale
        resolving it, and a bad portfolio name giving a readable error. They also check
        every other tool's numbers against the M1a tables (e.g. Tech net 37.0\%, benchmark
        40.0\%, active $-3.0$ on 10/21).
  \item One test launches \code{python -m pm_agent.mcp_server} as a subprocess and calls a
        tool over stdio, the way Claude Code will.
  \item Mutation check: with the \code{ToolError} conversion disabled, the 4 error tests
        failed; restored, all passed.
  \item Printed every tool description as the client receives it and read them as a model
        would. That is how the indentation problem was found.
  \item Full suite: 76 passed.
\end{enumerate}

\begin{reviewq}
\begin{enumerate}
  \item What four things does the server send for each tool, and which one does the model
        rely on most when choosing a tool?
  \item What would the model see if \code{check_limits} raised a plain \code{ValueError}
        for an unknown portfolio? Why does the SDK behave that way?
  \item Why does \code{simulate_trades} return \code{traded_positions} when the after-flags
        are already there?
  \item Why is \code{warn_at} not converted to percent, while \code{net_warn_buffer} is?
  \item Why does every tool return the \code{date} it used?
  \item What is the difference between the in-process test client and the stdio test, and
        why have both?
  \item The trades parameter is a \code{TypedDict}. What would the model see in the schema if
        it were a plain \code{list[dict]}?
\end{enumerate}
\end{reviewq}

\begin{takeaway}
The MCP server is a translation layer: it loads data, calls the tested analytics, and turns
decimals into the units a reader expects. It does no math of its own. Most of the work was
in what the model sees: descriptions with example questions, named fields, the date on every
answer, and error messages that say what is available. One output, the traded stock's
weight after a trade, was added specifically so the model never has to compute a number
itself.
\end{takeaway}

% ---------------------------------------------------------------------
\section{M2b: the server inside Claude Code}

\subsection{Goal}

Tasks 2.2 and 2.3: connect our MCP server to a real client, Claude Code, and ask it portfolio
questions in plain English. Then fix any tool description that made it pick the wrong tool or
ask for clarification. This is the first time a model uses our tools, so it is the first real
test of the descriptions written in M2a.

\subsection{What was done}

\begin{enumerate}
  \item Ran \code{claude mcp add --scope project pm-agent -- .venv/bin/python -m
        pm_agent.mcp_server}, which wrote \code{.mcp.json} (Part~I). The first attempt failed
        with ``command not found'' because of the shell PATH (Part~I); one line in
        \code{~/.zshrc} fixed it.
  \item Approved the server for this project (but not ``all future servers'', see Part~I).
  \item \code{claude mcp list} reported \code{pm-agent ... Connected}.
  \item Asked six questions, one at a time, and checked every number in every answer
        against the tool output.
\end{enumerate}

\subsection{The six questions}

\begin{center}
\begin{tabularx}{\textwidth}{@{}X c X@{}}
\toprule
\textbf{Question} & \textbf{Calls} & \textbf{What happened} \\
\midrule
Are we breaching any limits? & 1 & NVDA 15.2\% vs 15\% breach; nothing else. Correct. \\
Five biggest positions; most overweight? & 2 & Correct. Saw only the top 10 bets, but reasoned correctly that any other overweight is at most 0.6 points. \\
What changed overnight to cause the breach? & 2 & NVDA $+10\%$, weight 14\% $\to$ 15.2\%. It also checked 10/21 to show warning $\to$ breach. Correct. \\
Sell \$1M NVDA, buy \$1M UNH: still breaching? & 1 & One call with both trades. Breach resolved, no new flags, NVDA 13.23\%, UNH 5.15\%. Correct. \\
Our weight in FAKECO? & 1 & ``Not held, not in the benchmark.'' No invented number. \\
Which stock should I buy next? & 1 & Declined to pick, but then framed benchmark underweights as things to buy. Borderline. \\
\bottomrule
\end{tabularx}
\end{center}

\textbf{Every number matched a tool result.} To check, the same tool calls were rerun with a
short script and compared value by value: market values, active weights, sector active
weights, overnight changes, and the after-trade weights (including UNH's 7{,}880.32 shares).
None were invented.

Notice what the model did \emph{not} need to do. For ``sell \$1M NVDA'', it read NVDA's new
weight (13.23\%) straight from \code{traded_positions}, the field added in M2a for exactly this
reason. For ``what caused the breach'', it chose to call \code{check_limits} for the previous
day too, without being told, because the description says what the tool is for.

\subsection{Task 2.3: description changes}

\textbf{None.} The task is to rewrite descriptions that caused a wrong tool choice or a
request for clarification, and neither happened. Changing descriptions that already work
would be change for its own sake, and would make it harder to tell later (in M4, with real
measurements) what helped.

Four observations were logged in \code{docs/NOTES.md} for later milestones instead:

\begin{enumerate}
  \item \textbf{The advice question is a prompt problem, not a tool problem.} With no
        system prompt (that arrives in M3), Claude still leaned towards suggesting buys. SPEC
        rule 4 (``explain the effect of trades the user proposes'') belongs in the agent's
        prompt.
  \item \textbf{\code{top_n} ranks by absolute size}, so for ``most overweight'' the top 10
        can be mostly underweights. Fine on 25 stocks; may need a direction filter with a
        500-stock benchmark (M6). Evals will tell.
  \item \textbf{No tool sizes a trade to a target weight.} Claude offered to work out ``how
        much NVDA to sell to get under 15\%''. With dollar-only trades, that means guessing
        and retrying, or doing math. The morning brief (M5) needs this, so it is checked there.
  \item \textbf{Loose summaries.} ``About $-0.1$ points'' summarised several small values.
        Harmless to a reader, but the M4 made-up-number check needs a rule for it.
\end{enumerate}

\subsection{Decisions and why}

\begin{itemize}
  \item \textbf{Project scope (\code{.mcp.json}, committed)} rather than local scope, so the
        setup travels with the repo, and \textbf{a relative command path} so it works on
        any clone.
  \item \textbf{Approve this server only}, not all future project servers, because the repo
        is public.
  \item \textbf{Change nothing without evidence.} Six good answers are not proof the
        descriptions are perfect, but they are no evidence of a problem either. M4's evals
        are where description changes get measured.
\end{itemize}

\subsection{How it was verified}

\begin{enumerate}
  \item \code{claude mcp list} shows the server connected (task 2.2's first check).
  \item ``Are we breaching any limits?'' was answered correctly using the tools (its second
        check).
  \item All numbers in all six answers were checked against a rerun of the same tool calls.
\end{enumerate}

\begin{reviewq}
\begin{enumerate}
  \item What is the difference between local, project and user scope for an MCP server, and
        why is project scope right here?
  \item Why did \code{claude} give ``command not found'' although it was installed?
  \item Why not choose ``use this and all future MCP servers in this project''?
  \item The advice answer was borderline. Why is the fix in the system prompt (M3) and not
        in a tool description?
  \item Task 2.3 changed nothing. Why is that a reasonable outcome, and not a skipped task?
  \item Which field added in M2a let the model answer the \$1M trade question without doing
        any math?
\end{enumerate}
\end{reviewq}

\begin{takeaway}
The first real model to use the tools picked the right one every time, asked for no
clarification, and invented no numbers. The descriptions were left as they were, because
nothing pointed to a problem. The rough edges it exposed (advice, ranking direction, sizing a
trade to a target) were logged with the milestone that should deal with each one.
\end{takeaway}

% ---------------------------------------------------------------------
\section{M3: the agent}

\subsection{Goal}

Tasks 3.1 to 3.4: our own agent, a program that takes a question in plain English, lets Claude
choose and call our MCP tools, and returns the answer with a full record of what happened.
Hand-written, with no agent framework (SPEC~\S8).

\begin{lstlisting}[style=plain]
$ python -m pm_agent.agent "Are we breaching any limits?"
Yes, there's one breach as of 2026-10-22. NVDA is at 15.2% of NAV, which is just over
the 15.0% maximum position size. ...
[claude-opus-5-5 | 2 model calls | tools: check_limits | 7395 in / 233 out tokens | 4.8s]
\end{lstlisting}

\subsection{What was built}

\begin{center}
\begin{tabularx}{\textwidth}{@{}l X@{}}
\toprule
\textbf{File} & \textbf{What it does} \\
\midrule
\code{src/pm_agent/agent.py}   & the loop (3.1), run log (3.3), command-line entry point \\
\code{src/pm_agent/prompts.py} & the system prompt with the five SPEC~\S8 rules (3.2) \\
\code{tests/test_agent.py}     & 12 tests using a fake model: no API calls, no cost \\
\code{logs/runs.jsonl}         & one line per run (git-ignored; created on first run) \\
\code{pyproject.toml}          & adds the \code{anthropic} SDK (the planned M3 dependency) \\
\bottomrule
\end{tabularx}
\end{center}

\subsection{The loop (3.1)}

The heart of \code{run_agent}, slightly shortened:

\begin{lstlisting}[style=py]
async with Client(server or server_params()) as mcp:            # start the MCP server
    tools = [to_claude_tool(t) for t in (await mcp.list_tools()).tools]
    messages = [{"role": "user", "content": question}]
    request = {"model": log.model, "max_tokens": MAX_TOKENS, "tools": tools,
               "system": f"{system or SYSTEM_PROMPT}\n\n{mcp.instructions}",
               "output_config": {"effort": EFFORT}}

    while log.steps < max_steps:
        response = await create_message(**request, messages=messages)
        log.steps += 1                                         # (token counting omitted)
        messages.append({"role": "assistant", "content": response.content})

        if response.stop_reason != "tool_use":                 # an answer: we're done
            log.answer = final_text(response.content)
            break

        results = []
        for block in response.content:
            if block.type == "tool_use":
                result = await mcp.call_tool(block.name, block.input)
                results.append({"type": "tool_result", "tool_use_id": block.id,
                                "content": result_text(result),
                                "is_error": bool(result.is_error)})
        messages.append({"role": "user", "content": results})   # all results, one message
    else:
        log.stop_reason = "max_steps"                           # guard: never loop forever
\end{lstlisting}

Points worth noticing:

\begin{itemize}
  \item \textbf{Tool definitions are not written twice.} They come from the server's
        \code{list_tools()} and are renamed field-for-field (\code{to_claude_tool}), so a
        description improved in \code{mcp_server.py} reaches Claude Code and our agent alike.
        The same goes for the server's \code{instructions} (units, naming), appended to the
        system prompt.
  \item \textbf{The reply goes back unchanged.} \code{response.content} is appended as-is,
        including any \code{thinking} blocks. The model needs to see its own
        \code{tool_use} block next to the matching \code{tool_result}, and Claude Opus 5.5
        requires the history to be append-only.
  \item \textbf{Errors are data, not crashes.} A tool error (unknown portfolio, bad date) goes
        back as a \code{tool_result} with \code{is_error: true} and the message from
        \code{data.py}, so the model can fix its call. It is also recorded in the log.
  \item \textbf{\code{while ... else}.} In Python, the \code{else} of a loop runs only if the
        loop ended without \code{break}: here, only if the step limit ran out.
\end{itemize}

\subsection{One run, message by message (3.4)}\label{sec:m3-walkthrough}

This is a real, recorded run of ``Are we breaching any limits?'' (2026-10-24, Claude Opus
5.5, effort \code{medium}). The loop made two calls to the API.

\textbf{Before call 1.} The agent starts the MCP server as a subprocess, reads the 8 tool
definitions and the server's instructions, and builds the system prompt (1{,}932 characters:
our five rules plus the units explanation).

\textbf{Call 1: request.} One message:

\begin{lstlisting}[style=plain]
system:   <rules + server instructions>        tools: <8 definitions>
messages: [ user: "Are we breaching any limits?" ]
\end{lstlisting}

\textbf{Call 1: response.} \code{stop_reason: tool_use}. One block:

\begin{lstlisting}[style=plain]
tool_use  id=toolu_01MeGD...  name=check_limits  input={}
\end{lstlisting}

The model chose the tool from its description alone, and left both arguments at their defaults
(portfolio \code{sample_fund}, latest date). Usage: 3{,}571 input tokens, 27 output tokens.
Almost all the input is the system prompt and the tool definitions; the question itself is
a handful of tokens.

\textbf{Between the calls: our code runs the tool.} The loop sees \code{tool_use}, calls
\code{check_limits} through MCP (32 milliseconds), and gets the structured result.

\textbf{Call 2: request.} The whole conversation again, now three messages:

\begin{lstlisting}[style=plain]
messages: [
  user:      "Are we breaching any limits?"
  assistant: [tool_use id=toolu_01MeGD... name=check_limits input={}]
  user:      [tool_result tool_use_id=toolu_01MeGD... is_error=false
              content={"portfolio": "sample_fund", "date": "2026-10-22",
                       "breach_count": 1, "warning_count": 0, "limits": {...},
                       "flags": [{"rule": "max_position", "subject": "NVDA",
                                  "value_pct": 15.2, "limit_pct": 15.0,
                                  "status": "breach"}]}]
]
\end{lstlisting}

Note the roles: the tool result is sent as a \code{user} message, because it comes from our side
of the conversation. The \code{tool_use_id} ties it to the request it answers.

\textbf{Call 2: response.} \code{stop_reason: end_turn}. One \code{text} block:

\begin{lstlisting}[style=plain]
Yes, there's one breach as of 2026-10-22. NVDA is at 15.2% of NAV, which is just over
the 15.0% maximum position size.
There are no warnings, so nothing else is close to a limit. ...
If you have a trade in mind, such as trimming NVDA by a certain dollar amount, I can
test whether it clears the breach without creating new ones.
\end{lstlisting}

Usage: 3{,}824 input tokens, 206 output tokens. Every number in it (15.2, 15.0, the limits, the
date) is in the tool result. ``Just over'' is the words-not-numbers escape hatch from rule~1.

\textbf{After the loop.} The \code{end_turn} ends the loop. The run log records 2 model calls,
one tool call (\code{check_limits}, no arguments, 0.032~s, no error), 7{,}395 input and 233
output tokens, and 4.8~s, and is appended to \code{logs/runs.jsonl}.

\textbf{Cost of this run}: $7{,}395 \times \$4/10^6 + 233 \times \$20/10^6 \approx \$0.030 +
\$0.005 = \$0.035$. Most of it is input, and most of the input is the same system prompt and
tool definitions sent twice. That is the case for prompt caching, logged in
\code{docs/NOTES.md} for the eval work in M4.

\subsection{The system prompt (3.2)}

\begin{lstlisting}[style=plain]
You help a portfolio manager understand their fund: holdings, exposures, risk limits,
and the effect of trades they are considering. Answer using the portfolio tools, which
compute every number in tested code.

Rules:
1. Every number in your answer must appear in a tool result from this conversation.
   Don't estimate, and don't calculate new numbers yourself: no differences, sums,
   ratios, or re-rounding of tool numbers. For example, don't turn "15.2% vs a 15%
   limit" into "0.2 points over". If you want a number no tool result contains, call a
   tool that provides it, or describe it in words ("just over the limit").
2. Quote percentages and dollar amounts exactly as the tools return them, with units.
3. If a tool returns an error, read it ... Fix the call and try again.
4. Don't give investment advice ... offer to test a trade they are considering.
5. Be concise. Lead with the direct answer ... say which date the numbers are for.
\end{lstlisting}

\textbf{Before and after.} The first live run of the agent, with no system prompt, answered
correctly but said NVDA was ``0.2 percentage points over''. That is $15.2 - 15.0$, a
subtraction the model did itself: a direct break of the project's one rule. Three questions,
before and after the prompt (the first ``before'' is our agent with no system prompt; the
other two are from the M2b trial in Claude Code, which had none of our rules either):

\begin{center}
\begin{tabularx}{\textwidth}{@{}l X X@{}}
\toprule
\textbf{Question} & \textbf{Before} & \textbf{After} \\
\midrule
Breaching any limits? & ``0.2 percentage points over'' (computed) & quotes 15.2\% and 15.0\% only \\
Which stock to buy? & declined, then listed underweights ``to buy'' & declines; offers to test a proposed trade \\
What changed overnight? & ``about $-0.1$ pp'' (in no tool result) & ``about a tenth of a point or less'' \\
\bottomrule
\end{tabularx}
\end{center}

\textbf{Why the prompt is written this way.}

\begin{itemize}
  \item \textbf{A concrete example of the failure} (``0.2 points over'') teaches more than an
        abstract rule. Models follow examples closely.
  \item \textbf{An honest escape hatch.} Forbidding all derived numbers without an alternative
        would leave the model unable to say ``just over''. Allowing words keeps answers
        natural without inventing figures.
  \item \textbf{Reasons, not shouting.} Current models follow calm, explained instructions
        well; capital letters and ``NEVER'' tend to make them over-cautious.
  \item \textbf{Units live in the server, not here.} The agent appends the MCP server's
        instructions, so the explanation of \code{_pct} and \code{_usd} is written once.
\end{itemize}

\subsection{Run logging (3.3)}

Each run's \code{RunLog} (question, model, answer, stop reason, steps, each tool call with
arguments, duration and error, tokens in and out, latency, timestamp) is appended as one line
of \code{logs/runs.jsonl}. This is the raw material for the M4 metrics (accuracy, tool choice,
cost, latency). Tests redirect the log to a temporary file, so running the tests never writes
to the real log.

\subsection{Testing without the API}

\begin{lstlisting}[style=py]
fake = FakeModel(
    reply("tool_use", text("Let me check."), ask_tool("check_limits")),
    reply("end_turn", text("Yes: NVDA is at 15.2% vs a 15% limit.")),
)
log = run("Are we breaching any limits?", fake)
assert [(c.name, c.args, c.error) for c in log.tool_calls] == [("check_limits", {}, None)]
\end{lstlisting}

The fake returns scripted replies (built from the SDK's own block types, so they look exactly
like real ones) and records every request. The tools are \emph{real}: the loop talks to our
actual MCP server, in-process. The tests cover the single tool call, the tool definitions and
server instructions reaching the model, a tool error going back as \code{is_error}, two tool
calls answered in one message, the step limit, the model chosen from \code{PM_AGENT_MODEL},
the \code{.env} loader, and the run log.

\subsection{Decisions and why}

\begin{enumerate}
  \item \textbf{A hand-written loop rather than the SDK's tool runner.} The spec asks for it:
        every line of the harness should be explainable, and the evals need hooks into each
        step. The loop is about 40 lines.
  \item \textbf{The model is an argument (\code{create_message}).} That one decision makes the
        loop testable for free, and is how evals will run it later.
  \item \textbf{Effort set explicitly to \code{medium}.} It is the model's default, but
        writing it down makes the cost and quality lever visible, ready for M4 to compare.
  \item \textbf{Server-side refusal fallbacks enabled} (\code{fallbacks="default"}, a beta).
        If a classifier wrongly declines a question, the API retries on another model. The
        API guidance recommends this by default for this model. Removing it is two lines.
  \item \textbf{A ten-line \code{.env} loader instead of the \code{python-dotenv}
        library}, to avoid a dependency for something this small.
  \item \textbf{The step limit counts model calls} (default 10). Each call can request
        several tools, so this bounds cost per question.
  \item \textbf{No retries of our own.} The SDK already retries rate limits and server errors
        twice with backoff.
\end{enumerate}

\subsection{How it was verified}

\begin{enumerate}
  \item Wrote the fake-model tests first and watched them fail (no \code{agent} module, then
        no \code{prompts} module, then no \code{RUNS_LOG}); implemented until they passed.
  \item Ran \code{python -m pm_agent.agent "Are we breaching any limits?"} against the real
        API: correct answer, one \code{check_limits} call (task 3.1's ``done when'').
  \item Ran three questions with and without the system prompt and checked every number
        against the tool outputs (the table above).
  \item Recorded one full run's requests and responses for the walkthrough, and checked that
        the run appeared in \code{logs/runs.jsonl}.
  \item Full suite: 88 passed. Total API spend for M3: about \$0.20.
\end{enumerate}

\begin{reviewq}
\begin{enumerate}
  \item What does \code{stop_reason} tell the loop, and what does the loop do for
        \code{tool_use} versus \code{end_turn}?
  \item Why is the tool result sent with role \code{user}? What links it to the request?
  \item The API is stateless. What does that mean for what each request contains, and for
        cost?
  \item In the walkthrough, call 1 used 3{,}571 input tokens for a four-word question.
        Where did they come from?
  \item Why does the agent take \code{create_message} as an argument?
  \item Before the system prompt, the agent said ``0.2 percentage points over''. Why is that a
        problem when 15.2 and 15.0 both came from a tool?
  \item What does the \code{else} on the \code{while} loop do?
  \item Why are tool errors sent back to the model instead of raising an exception?
\end{enumerate}
\end{reviewq}

\begin{takeaway}
The agent is a short loop around one field, \code{stop_reason}: while the model asks for
tools, run them and send the results back; when it answers, stop. Everything else is
bookkeeping (the run log) and guard rails (the step limit, the system prompt). The first live
run showed why the guard rails matter: with no prompt, the model did a small subtraction on its
own. A system prompt with a concrete example and an honest escape hatch fixed it, and the
fake-model tests mean the loop itself is checked on every push without spending anything.
\end{takeaway}

% =====================================================================
\appendix
\section{Glossary}

\begin{description}[style=sameline,leftmargin=3.2cm]
  \item[Active weight] Portfolio weight minus benchmark weight.
  \item[Agent] An LLM in a loop that can call tools until it reaches an answer.
  \item[Benchmark] The index a fund is compared against.
  \item[Breach] A limit exceeded.
  \item[CI] Continuous integration: automatic tests on every push.
  \item[Context manager] An object used with \code{with} that runs setup and cleanup around a block.
  \item[Dataclass] A Python class that just holds typed fields; \code{frozen} makes it immutable.
  \item[DataFrame] A pandas table with named columns.
  \item[Decorator] A function applied with \code{@name} that registers or wraps another function.
  \item[Effort] How much the model reasons before answering (low to max); trades cost for depth.
  \item[Eval] A scored test suite for AI behaviour.
  \item[Fake (test double)] A scripted stand-in for something costly, used in tests.
  \item[Flag identity] A flag's (rule, subject) pair, used to compare flags before and after a trade.
  \item[Gross exposure] Long minus short: the total size of all bets.
  \item[GICS] The standard scheme for classifying companies into sectors.
  \item[Harness] The program around the model: loop, logging, limits.
  \item[Hedge] A position taken to offset risk from another.
  \item[Immutable] Cannot be changed after creation; changes make a new copy.
  \item[Index alignment] pandas matching rows by label when combining tables.
  \item[JSON Schema] A standard JSON description of a tool's arguments and their types.
  \item[LLM] Large language model.
  \item[.mcp.json] The committed file listing a project's MCP servers for Claude Code.
  \item[JSON Lines] A file with one JSON object per line; easy to append to.
  \item[MCP] Model Context Protocol: a standard for exposing tools to AI.
  \item[Messages API] Anthropic's endpoint for conversations with Claude.
  \item[Mutation check] Breaking code or data on purpose to prove a test can fail.
  \item[NAV] Net asset value: market values plus cash.
  \item[Net exposure] Long plus short: the direction of the market bet.
  \item[PATH] The list of folders the shell searches for commands.
  \item[Pre-trade check] Simulating a trade and checking limits before placing it.
  \item[Pure function] Output depends only on inputs; no files, network, or side effects.
  \item[Refusal fallback] Retrying a declined request on another model, server-side.
  \item[Seed] The starting value for a random number generator.
  \item[Server instructions] Guidance an MCP server sends once, for every tool.
  \item[Short] Selling borrowed shares; profits when the price falls.
  \item[Stable sort] A sort that keeps tied items in their original order.
  \item[Stop reason] Why the model stopped: \code{tool_use}, \code{end_turn}, ...
  \item[Structured output] A tool result as a JSON object with named fields.
  \item[System prompt] Standing instructions sent with every request.
  \item[Tool annotation] A behaviour hint on a tool, e.g. read-only.
  \item[Tool use] A model asking the program to run a function.
  \item[tool\_result / tool\_use] The content blocks that carry a tool's output and a tool request.
  \item[ToolError] The MCP exception whose message is shown to the model.
  \item[TypedDict] A type hint for a dict with known keys.
  \item[Warning] Close to a limit (90\% of it, or within 2 points for net).
  \item[Weight] Market value divided by NAV.
\end{description}

\section{Command cheat sheet}

\begin{lstlisting}[style=plain]
python3 -m venv .venv               # create a virtual environment
source .venv/bin/activate           # activate it (prompt shows (.venv))
pip install -e ".[dev]"             # install the project + dev tools
pytest -q                           # run all tests quietly
pytest tests/test_data.py -v        # run one file, verbose
python scripts/make_sample_data.py  # regenerate data/sample
python -m pm_agent.mcp_server       # run the MCP server over stdio
python -m pm_agent.agent "..."      # ask the agent a question (uses the API)

git status                          # what changed?
git diff                            # show the changes
git add -A && git commit -m "msg"   # snapshot everything
git push                            # upload to GitHub
git log --oneline                   # history, one line per commit
gh run list                         # recent GitHub Actions runs
\end{lstlisting}

\end{document}
